\documentclass[../main.tex]{subfiles}
\begin{document}

\chapter{Distributed File Systems}
\lessoninfo{
  video={P4L2},
  textbook={OSTEP \cite{ostep} ch.~48--49; Silberschatz et al.\ \cite{silberschatz2018} ch.~15, 19; Nelson et al.\ \cite{nelson1988}},
  keywords={distributed file system, replication, partitioning, upload/download model, stateless server, stateful server, UNIX semantics, session semantics, lease, NFS, file handle, Sprite, delayed write, concurrent write sharing},
}

The virtual file system of Lesson~11 hides from applications where the
files they use reside, and the remote procedure calls of Lesson~13 allow a
process to have operations performed on another machine.  A distributed
file system combines the two: it stores files on machines connected by a
network and serves them to clients that use the ordinary file system
interface.  This lesson examines the design space of such systems---how
the work is divided between clients and servers, what state a server
keeps, how cached copies are kept consistent and what a client observes
when files are shared---and then studies two concrete designs: NFS, and
Sprite, whose design was derived from measurements of how files are used.

\section{Distributed file systems}
\label{sec:l14-dfs}

Operating systems hide the details of storage behind the file system
interface.\source{P4L2, lesson preview, visual metaphor, distributed file
systems, DFS models; OSTEP ch.~48; Silberschatz et al.\ ch.~19.}  Beneath
it, the virtual file system passes each operation to the file system in
which the file resides, and that file system may keep its files on a local
device or on other machines of a network.

\begin{definition}[Distributed file system]
  A \term{distributed file system}[分散ファイルシステム]%
  \index{DFS|see{distributed file system}} (DFS) is a file system whose
  files are stored on, and served by, machines connected by a network, and
  which processes access through the ordinary file system interface
  whether or not a file resides on their own machine.
\end{definition}

Three properties characterize a DFS:
\begin{itemize}
  \item It is accessed through a well-defined interface.  Applications use
    the same system calls as for local files, and the VFS passes each
    operation on a remote file to the DFS.
  \item It must maintain a consistent state.  When several clients use the
    same files, the DFS must track which client accesses which file,
    propagate updates and keep cached copies coherent.
  \item It may mix several distribution models.  Files can be replicated or
    partitioned among several servers, and the roles of client and server
    need not reside on different machines.
\end{itemize}

The distribution models are shown in \cref{fig:l14-models}.
\begin{description}
  \item[Client and server on different machines] A single file server
    stores all files, and clients on other machines access them over the
    network.
  \item[Distributed file server] The file service itself runs on several
    machines.  In a \emph{replicated} file server every machine holds all
    files; in a \emph{partitioned} file server each machine holds a part of
    the files.  \Cref{sec:l14-replication} compares the two.
  \item[Peers] Files are stored on and served from all machines.
\end{description}

In the last model the distinction between clients and servers is blurred.
Such an organization is called \term{peer-to-peer}[ピアツーピア]: every
machine is a peer that both uses and provides the file service.  Each peer
stores a portion of the files and carries part of the load by serving
requests, typically those for the files stored on it.

\begin{figure}[htbp]
  \centering
  % DFS models: (a) client and server on different machines, (b) replicated
% file servers, (c) partitioned file servers, (d) peers.
% Usage: % DFS models: (a) client and server on different machines, (b) replicated
% file servers, (c) partitioned file servers, (d) peers.
% Usage: % DFS models: (a) client and server on different machines, (b) replicated
% file servers, (c) partitioned file servers, (d) peers.
% Usage: \input{figures/l14-models}   (width ~118 mm)
\begin{tikzpicture}[x=1mm, y=1mm, font=\scriptsize\sffamily,
  cl/.style={draw, fill=white, minimum width=12mm, minimum height=5mm, inner sep=1pt},
  sv/.style={draw, fill=ln-box, minimum width=15mm, minimum height=8mm, inner sep=1pt, align=center},
  pe/.style={draw, fill=ln-accent!12, minimum width=15mm, minimum height=8mm, inner sep=1pt, align=center},
  net/.style={line width=1.2pt, ln-muted},
  lab/.style={font=\scriptsize\sffamily, text=ln-muted},
  ttl/.style={font=\footnotesize\sffamily\bfseries, anchor=west}]
  % (a) one server
  \begin{scope}[shift={(0,33)}]
    \node[ttl] at (0,27) {\iflangja{(a) 単一のファイルサーバ}{(a) single file server}};
    \draw[net] (3,13) -- (53,13);
    \foreach \x in {10,28,46} {
      \node[cl] at (\x,20) {\iflangja{クライアント}{client}};
      \draw (\x,17.5) -- (\x,13);
    }
    \node[lab, anchor=south] at (19,13.3) {\iflangja{ネットワーク}{network}};
    \node[sv] at (28,4) {\iflangja{サーバ}{server}\\A--F};
    \draw (28,8) -- (28,13);
  \end{scope}
  % (b) replicated
  \begin{scope}[shift={(62,33)}]
    \node[ttl] at (0,27) {\iflangja{(b) 複製されたファイルサーバ}{(b) replicated file servers}};
    \draw[net] (3,13) -- (53,13);
    \foreach \x in {10,28,46} {
      \node[cl] at (\x,20) {\iflangja{クライアント}{client}};
      \draw (\x,17.5) -- (\x,13);
      \node[sv] at (\x,4) {\iflangja{サーバ}{server}\\A--F};
      \draw (\x,8) -- (\x,13);
    }
  \end{scope}
  % (c) partitioned
  \begin{scope}[shift={(0,0)}]
    \node[ttl] at (0,27) {\iflangja{(c) 分割されたファイルサーバ}{(c) partitioned file servers}};
    \draw[net] (3,13) -- (53,13);
    \foreach \x/\f in {10/{A, B},28/{C, D},46/{E, F}} {
      \node[cl] at (\x,20) {\iflangja{クライアント}{client}};
      \draw (\x,17.5) -- (\x,13);
      \node[sv] at (\x,4) {\iflangja{サーバ}{server}\\\f};
      \draw (\x,8) -- (\x,13);
    }
  \end{scope}
  % (d) peers
  \begin{scope}[shift={(62,0)}]
    \node[ttl] at (0,27) {\iflangja{(d) ピア}{(d) peers}};
    \draw[net] (3,19) -- (53,19);
    \foreach \x/\f in {10/{A, B},28/{C, D},46/{E, F}} {
      \node[pe] at (\x,10) {\iflangja{ピア}{peer}\\\f};
      \draw (\x,14) -- (\x,19);
    }
    \node[lab, anchor=north] at (28,5)
      {\iflangja{各ピアがクライアントでありサーバでもある}{each peer is both client and server}};
  \end{scope}
\end{tikzpicture}
   (width ~118 mm)
\begin{tikzpicture}[x=1mm, y=1mm, font=\scriptsize\sffamily,
  cl/.style={draw, fill=white, minimum width=12mm, minimum height=5mm, inner sep=1pt},
  sv/.style={draw, fill=ln-box, minimum width=15mm, minimum height=8mm, inner sep=1pt, align=center},
  pe/.style={draw, fill=ln-accent!12, minimum width=15mm, minimum height=8mm, inner sep=1pt, align=center},
  net/.style={line width=1.2pt, ln-muted},
  lab/.style={font=\scriptsize\sffamily, text=ln-muted},
  ttl/.style={font=\footnotesize\sffamily\bfseries, anchor=west}]
  % (a) one server
  \begin{scope}[shift={(0,33)}]
    \node[ttl] at (0,27) {\iflangja{(a) 単一のファイルサーバ}{(a) single file server}};
    \draw[net] (3,13) -- (53,13);
    \foreach \x in {10,28,46} {
      \node[cl] at (\x,20) {\iflangja{クライアント}{client}};
      \draw (\x,17.5) -- (\x,13);
    }
    \node[lab, anchor=south] at (19,13.3) {\iflangja{ネットワーク}{network}};
    \node[sv] at (28,4) {\iflangja{サーバ}{server}\\A--F};
    \draw (28,8) -- (28,13);
  \end{scope}
  % (b) replicated
  \begin{scope}[shift={(62,33)}]
    \node[ttl] at (0,27) {\iflangja{(b) 複製されたファイルサーバ}{(b) replicated file servers}};
    \draw[net] (3,13) -- (53,13);
    \foreach \x in {10,28,46} {
      \node[cl] at (\x,20) {\iflangja{クライアント}{client}};
      \draw (\x,17.5) -- (\x,13);
      \node[sv] at (\x,4) {\iflangja{サーバ}{server}\\A--F};
      \draw (\x,8) -- (\x,13);
    }
  \end{scope}
  % (c) partitioned
  \begin{scope}[shift={(0,0)}]
    \node[ttl] at (0,27) {\iflangja{(c) 分割されたファイルサーバ}{(c) partitioned file servers}};
    \draw[net] (3,13) -- (53,13);
    \foreach \x/\f in {10/{A, B},28/{C, D},46/{E, F}} {
      \node[cl] at (\x,20) {\iflangja{クライアント}{client}};
      \draw (\x,17.5) -- (\x,13);
      \node[sv] at (\x,4) {\iflangja{サーバ}{server}\\\f};
      \draw (\x,8) -- (\x,13);
    }
  \end{scope}
  % (d) peers
  \begin{scope}[shift={(62,0)}]
    \node[ttl] at (0,27) {\iflangja{(d) ピア}{(d) peers}};
    \draw[net] (3,19) -- (53,19);
    \foreach \x/\f in {10/{A, B},28/{C, D},46/{E, F}} {
      \node[pe] at (\x,10) {\iflangja{ピア}{peer}\\\f};
      \draw (\x,14) -- (\x,19);
    }
    \node[lab, anchor=north] at (28,5)
      {\iflangja{各ピアがクライアントでありサーバでもある}{each peer is both client and server}};
  \end{scope}
\end{tikzpicture}
   (width ~118 mm)
\begin{tikzpicture}[x=1mm, y=1mm, font=\scriptsize\sffamily,
  cl/.style={draw, fill=white, minimum width=12mm, minimum height=5mm, inner sep=1pt},
  sv/.style={draw, fill=ln-box, minimum width=15mm, minimum height=8mm, inner sep=1pt, align=center},
  pe/.style={draw, fill=ln-accent!12, minimum width=15mm, minimum height=8mm, inner sep=1pt, align=center},
  net/.style={line width=1.2pt, ln-muted},
  lab/.style={font=\scriptsize\sffamily, text=ln-muted},
  ttl/.style={font=\footnotesize\sffamily\bfseries, anchor=west}]
  % (a) one server
  \begin{scope}[shift={(0,33)}]
    \node[ttl] at (0,27) {\iflangja{(a) 単一のファイルサーバ}{(a) single file server}};
    \draw[net] (3,13) -- (53,13);
    \foreach \x in {10,28,46} {
      \node[cl] at (\x,20) {\iflangja{クライアント}{client}};
      \draw (\x,17.5) -- (\x,13);
    }
    \node[lab, anchor=south] at (19,13.3) {\iflangja{ネットワーク}{network}};
    \node[sv] at (28,4) {\iflangja{サーバ}{server}\\A--F};
    \draw (28,8) -- (28,13);
  \end{scope}
  % (b) replicated
  \begin{scope}[shift={(62,33)}]
    \node[ttl] at (0,27) {\iflangja{(b) 複製されたファイルサーバ}{(b) replicated file servers}};
    \draw[net] (3,13) -- (53,13);
    \foreach \x in {10,28,46} {
      \node[cl] at (\x,20) {\iflangja{クライアント}{client}};
      \draw (\x,17.5) -- (\x,13);
      \node[sv] at (\x,4) {\iflangja{サーバ}{server}\\A--F};
      \draw (\x,8) -- (\x,13);
    }
  \end{scope}
  % (c) partitioned
  \begin{scope}[shift={(0,0)}]
    \node[ttl] at (0,27) {\iflangja{(c) 分割されたファイルサーバ}{(c) partitioned file servers}};
    \draw[net] (3,13) -- (53,13);
    \foreach \x/\f in {10/{A, B},28/{C, D},46/{E, F}} {
      \node[cl] at (\x,20) {\iflangja{クライアント}{client}};
      \draw (\x,17.5) -- (\x,13);
      \node[sv] at (\x,4) {\iflangja{サーバ}{server}\\\f};
      \draw (\x,8) -- (\x,13);
    }
  \end{scope}
  % (d) peers
  \begin{scope}[shift={(62,0)}]
    \node[ttl] at (0,27) {\iflangja{(d) ピア}{(d) peers}};
    \draw[net] (3,19) -- (53,19);
    \foreach \x/\f in {10/{A, B},28/{C, D},46/{E, F}} {
      \node[pe] at (\x,10) {\iflangja{ピア}{peer}\\\f};
      \draw (\x,14) -- (\x,19);
    }
    \node[lab, anchor=north] at (28,5)
      {\iflangja{各ピアがクライアントでありサーバでもある}{each peer is both client and server}};
  \end{scope}
\end{tikzpicture}

  \caption{DFS models.  (a) A single server stores all files.  (b) Every
    server of a replicated file server holds all files.  (c) Each server of
    a partitioned file server holds some of the files.  (d) All machines
    are peers that store files and serve them.}
  \label{fig:l14-models}
\end{figure}

\section{Remote file service}
\label{sec:l14-remote}

Even with a single server, the work of accessing a remote file can be
divided between client and server in very different ways.\source{P4L2,
remote file service: extremes, remote file service: a compromise; OSTEP
ch.~49; Silberschatz et al.\ ch.~19.}  Two extremes bound the design space
(\cref{fig:l14-remote}).

\begin{definition}[Upload/download model]
  In the \term{upload/download model}[アップロード・ダウンロードモデル] a
  client that accesses a file downloads the entire file from the server,
  performs all operations on its local copy, and uploads the entire file to
  the server when it is done.
\end{definition}

File transfer with FTP and version control systems such as Subversion work
in this way.  Reads and writes are performed locally at the client and
are therefore fast.  But the entire file is transferred even when the
client needs only a small part of it, and the server gives up control: once
a file has been downloaded, the server knows neither what the client does
with it nor when the modified file will be uploaded.

\begin{definition}[Remote access model]
  In the \term{remote access model}[遠隔アクセスモデル], also called true
  remote file access, the file remains on the server and every operation on
  it---every read and every write---is sent to the server.  Nothing is done
  locally at the client.
\end{definition}

Because every access passes through the server, file access is
centralized, and consistency is as easy to reason about as for a local
file system shared by several processes.  But every file operation pays
for a network round trip and for processing at the server, and since the
server executes every operation of every client, its capacity limits the
number of clients the DFS can support.

\begin{figure}[htbp]
  \centering
  % Remote file service: (a) upload/download model, (b) remote access model,
% (c) the compromise in which clients cache some blocks.
% Usage: % Remote file service: (a) upload/download model, (b) remote access model,
% (c) the compromise in which clients cache some blocks.
% Usage: % Remote file service: (a) upload/download model, (b) remote access model,
% (c) the compromise in which clients cache some blocks.
% Usage: \input{figures/l14-remote}   (width ~118 mm)
\begin{tikzpicture}[x=1mm, y=1mm, font=\scriptsize\sffamily,
  >={Stealth[length=1.5mm]},
  machine/.style={draw, fill=white, rounded corners=0.8mm},
  lab/.style={font=\scriptsize\sffamily, align=center},
  who/.style={font=\scriptsize\sffamily\itshape, text=ln-muted, anchor=north west},
  ttl/.style={font=\footnotesize\sffamily\bfseries, anchor=west}]
  % \lblocks{x0}{y0}{list of block indices to fill}{fill colour}
  \def\lblocks#1#2#3#4{%
    \foreach \i in {0,...,5} { \draw[fill=white] (#1+\i*4,#2) rectangle ++(4,4); }
    \foreach \i in {#3} { \draw[fill=#4] (#1+\i*4,#2) rectangle ++(4,4); }}
  \foreach \px/\t in {0/{\iflangja{(a) アップロード・ダウンロード}{(a) upload/download}},
                      40/{\iflangja{(b) 遠隔アクセス}{(b) remote access}},
                      80/{\iflangja{(c) 折衷案}{(c) compromise}}} {
    \node[ttl] at (\px,44) {\t};
    \draw[machine] (\px+1,26) rectangle (\px+37,39);
    \node[who] at (\px+1.5,38.6) {\iflangja{クライアント}{client}};
    \draw[machine, fill=ln-box] (\px+1,0) rectangle (\px+37,11);
    \node[who] at (\px+1.5,10.6) {\iflangja{サーバ}{server}};
    \lblocks{\px+7}{1.8}{0,...,5}{ln-muted!45}
  }
  % (a) whole file at the client
  \lblocks{7}{28}{0,...,5}{ln-accent!45}
  \draw[->, thick] (9,11.5) -- node[lab, right, pos=0.7] {\iflangja{ダウンロード}{download}} (9,25.5);
  \draw[->, thick] (29,25.5) -- node[lab, left, pos=0.7] {\iflangja{アップロード}{upload}} (29,11.5);
  % (b) nothing at the client, every operation to the server
  \node[lab, text=ln-muted] at (59,30) {\iflangja{ローカルにはデータなし}{no file data kept locally}};
  \foreach \x in {50,54,58} { \draw[->] (\x,25.5) -- (\x,11.5); }
  \foreach \x in {62,66,70} { \draw[<-] (\x,25.5) -- (\x,11.5); }
  \node[lab, fill=white, inner sep=0.6pt] at (60,18.5) {\iflangja{すべての操作}{every operation}};
  % (c) some blocks cached, frequent interaction
  \lblocks{87}{28}{1,2}{ln-accent!45}
  \draw[->, thick] (89,11.5) -- node[lab, right, pos=0.25] {\iflangja{ブロック}{blocks}} (89,25.5);
  \draw[->, thick] (109,25.5) -- node[lab, left, pos=0.25] {\iflangja{open，検証，\\書き戻し}{open, check,\\write back}} (109,11.5);
\end{tikzpicture}
   (width ~118 mm)
\begin{tikzpicture}[x=1mm, y=1mm, font=\scriptsize\sffamily,
  >={Stealth[length=1.5mm]},
  machine/.style={draw, fill=white, rounded corners=0.8mm},
  lab/.style={font=\scriptsize\sffamily, align=center},
  who/.style={font=\scriptsize\sffamily\itshape, text=ln-muted, anchor=north west},
  ttl/.style={font=\footnotesize\sffamily\bfseries, anchor=west}]
  % \lblocks{x0}{y0}{list of block indices to fill}{fill colour}
  \def\lblocks#1#2#3#4{%
    \foreach \i in {0,...,5} { \draw[fill=white] (#1+\i*4,#2) rectangle ++(4,4); }
    \foreach \i in {#3} { \draw[fill=#4] (#1+\i*4,#2) rectangle ++(4,4); }}
  \foreach \px/\t in {0/{\iflangja{(a) アップロード・ダウンロード}{(a) upload/download}},
                      40/{\iflangja{(b) 遠隔アクセス}{(b) remote access}},
                      80/{\iflangja{(c) 折衷案}{(c) compromise}}} {
    \node[ttl] at (\px,44) {\t};
    \draw[machine] (\px+1,26) rectangle (\px+37,39);
    \node[who] at (\px+1.5,38.6) {\iflangja{クライアント}{client}};
    \draw[machine, fill=ln-box] (\px+1,0) rectangle (\px+37,11);
    \node[who] at (\px+1.5,10.6) {\iflangja{サーバ}{server}};
    \lblocks{\px+7}{1.8}{0,...,5}{ln-muted!45}
  }
  % (a) whole file at the client
  \lblocks{7}{28}{0,...,5}{ln-accent!45}
  \draw[->, thick] (9,11.5) -- node[lab, right, pos=0.7] {\iflangja{ダウンロード}{download}} (9,25.5);
  \draw[->, thick] (29,25.5) -- node[lab, left, pos=0.7] {\iflangja{アップロード}{upload}} (29,11.5);
  % (b) nothing at the client, every operation to the server
  \node[lab, text=ln-muted] at (59,30) {\iflangja{ローカルにはデータなし}{no file data kept locally}};
  \foreach \x in {50,54,58} { \draw[->] (\x,25.5) -- (\x,11.5); }
  \foreach \x in {62,66,70} { \draw[<-] (\x,25.5) -- (\x,11.5); }
  \node[lab, fill=white, inner sep=0.6pt] at (60,18.5) {\iflangja{すべての操作}{every operation}};
  % (c) some blocks cached, frequent interaction
  \lblocks{87}{28}{1,2}{ln-accent!45}
  \draw[->, thick] (89,11.5) -- node[lab, right, pos=0.25] {\iflangja{ブロック}{blocks}} (89,25.5);
  \draw[->, thick] (109,25.5) -- node[lab, left, pos=0.25] {\iflangja{open，検証，\\書き戻し}{open, check,\\write back}} (109,11.5);
\end{tikzpicture}
   (width ~118 mm)
\begin{tikzpicture}[x=1mm, y=1mm, font=\scriptsize\sffamily,
  >={Stealth[length=1.5mm]},
  machine/.style={draw, fill=white, rounded corners=0.8mm},
  lab/.style={font=\scriptsize\sffamily, align=center},
  who/.style={font=\scriptsize\sffamily\itshape, text=ln-muted, anchor=north west},
  ttl/.style={font=\footnotesize\sffamily\bfseries, anchor=west}]
  % \lblocks{x0}{y0}{list of block indices to fill}{fill colour}
  \def\lblocks#1#2#3#4{%
    \foreach \i in {0,...,5} { \draw[fill=white] (#1+\i*4,#2) rectangle ++(4,4); }
    \foreach \i in {#3} { \draw[fill=#4] (#1+\i*4,#2) rectangle ++(4,4); }}
  \foreach \px/\t in {0/{\iflangja{(a) アップロード・ダウンロード}{(a) upload/download}},
                      40/{\iflangja{(b) 遠隔アクセス}{(b) remote access}},
                      80/{\iflangja{(c) 折衷案}{(c) compromise}}} {
    \node[ttl] at (\px,44) {\t};
    \draw[machine] (\px+1,26) rectangle (\px+37,39);
    \node[who] at (\px+1.5,38.6) {\iflangja{クライアント}{client}};
    \draw[machine, fill=ln-box] (\px+1,0) rectangle (\px+37,11);
    \node[who] at (\px+1.5,10.6) {\iflangja{サーバ}{server}};
    \lblocks{\px+7}{1.8}{0,...,5}{ln-muted!45}
  }
  % (a) whole file at the client
  \lblocks{7}{28}{0,...,5}{ln-accent!45}
  \draw[->, thick] (9,11.5) -- node[lab, right, pos=0.7] {\iflangja{ダウンロード}{download}} (9,25.5);
  \draw[->, thick] (29,25.5) -- node[lab, left, pos=0.7] {\iflangja{アップロード}{upload}} (29,11.5);
  % (b) nothing at the client, every operation to the server
  \node[lab, text=ln-muted] at (59,30) {\iflangja{ローカルにはデータなし}{no file data kept locally}};
  \foreach \x in {50,54,58} { \draw[->] (\x,25.5) -- (\x,11.5); }
  \foreach \x in {62,66,70} { \draw[<-] (\x,25.5) -- (\x,11.5); }
  \node[lab, fill=white, inner sep=0.6pt] at (60,18.5) {\iflangja{すべての操作}{every operation}};
  % (c) some blocks cached, frequent interaction
  \lblocks{87}{28}{1,2}{ln-accent!45}
  \draw[->, thick] (89,11.5) -- node[lab, right, pos=0.25] {\iflangja{ブロック}{blocks}} (89,25.5);
  \draw[->, thick] (109,25.5) -- node[lab, left, pos=0.25] {\iflangja{open，検証，\\書き戻し}{open, check,\\write back}} (109,11.5);
\end{tikzpicture}

  \caption{Remote file service.  (a) The client downloads the whole file
    and uploads it when done.  (b) The file stays at the server, which
    performs every operation.  (c) The client caches some blocks but
    interacts with the server frequently.}
  \label{fig:l14-remote}
\end{figure}

\begin{example}
  A client works with a \SI{1}{\gibi\byte} file over a network that
  transfers \SI[per-mode=symbol]{100}{\mebi\byte\per\second}; every request
  costs a round trip of \SI{1}{\milli\second}.  If the client reads 50
  blocks of \SI{4}{\kibi\byte}, the remote access model needs 50 requests,
  about \SI{50}{\milli\second}, whereas downloading the file takes
  $1024/100 \approx \SI{10.2}{\second}$.  If the client instead reads
  \num{200000} blocks, the round trips of the remote access model alone
  take \SI{200}{\second}, while the download still takes
  \SI{10.2}{\second} and every read after it is local.
\end{example}

\subsection{A compromise}

Practical DFSs lie between the extremes and combine two principles:
\begin{enumerate}
  \item Clients may keep parts of files, individual blocks, locally.
    Operations on local blocks have low latency, and the requests they
    save reduce the load on the server, so that a server can support more
    clients.
  \item Clients must interact with the server, and do so frequently: when
    they open a file, to check whether their local blocks are still
    current, and to send their modifications back.  The server thereby
    gains insight into what the clients are doing and control over which
    accesses it permits, which makes consistency easier to maintain.
\end{enumerate}
The price is a more complex server, and file sharing semantics that differ
from those of a local file system: a client that works on local copies of
blocks cannot observe every change made by another client at the moment
it happens (\cref{sec:l14-semantics}).

\section{Stateless and stateful file servers}
\label{sec:l14-state}

The division of work between client and server determines what the
server must remember about its clients.\source{P4L2, stateless vs.\
stateful file server; OSTEP ch.~49; Silberschatz et al.\ ch.~19.}
Conversely, the state that a server keeps determines which of the models
of \cref{sec:l14-remote} it can support.

\begin{definition}[Stateless server]
  A \term{stateless server}[ステートレスサーバ] keeps no information about
  its clients between requests: neither which files are open, nor which
  data a client has cached, nor where in a file a client is reading.  Every
  request is self-contained and carries all information needed to serve
  it.
\end{definition}

\begin{definition}[Stateful server]
  A \term{stateful server}[ステートフルサーバ] maintains information about
  its clients: which clients access which files, in which mode, and which
  data they have cached or locked.
\end{definition}

A stateless server suffices for the two extreme models but not for the
compromise, because keeping the blocks cached by clients consistent
requires knowing which clients have cached which files.  Its advantages
are that it spends no memory or CPU time on client state, and that it
recovers from a failure simply by restarting: no state has been lost, and
a client whose request went unanswered sends it again.  Resending is safe
when requests are \emph{idempotent}, that is, when executing a request
twice has the same effect as executing it once, as for a read or a write at
an explicit offset.  Its disadvantages are that it cannot support caching
with consistency management, nor locking, and that every request must
describe the operation completely---which file, which offset, how many
bytes---so that more bits are transferred per request.

A stateful server is required for the compromise model.  Because it knows
what its clients do, it can support locking, client caching with
consistency guarantees, and incremental operations such as reading the
next block of a file.  The price is twofold.  When the server fails, its
state is lost unless the server checkpoints it (Lesson~8)%
\index{checkpointing} and a recovery mechanism restores it; and
maintaining the state and enforcing consistency costs server resources,
how much depending on the caching mechanism and the consistency protocol.
\Cref{tab:l14-state} summarizes the comparison.

\begin{table}[htbp]
  \centering
  \caption{Stateless and stateful file servers.}
  \label{tab:l14-state}
  \small
  \begin{tabular}{@{}>{\raggedright\arraybackslash}p{0.2\linewidth}>{\raggedright\arraybackslash}p{0.36\linewidth}>{\raggedright\arraybackslash}p{0.36\linewidth}@{}}
    \toprule
    & Stateless server & Stateful server \\
    \midrule
    Suited to & upload/download and remote access models & the compromise
      model \\
    Requests & self-contained; more bits per request & may refer to state
      kept by the server \\
    Supports & neither consistent caching nor locking & locking, caching,
      incremental operations \\
    Server resources & none for client state & memory and CPU time for state
      and consistency \\
    After a failure & restart and continue & state must be checkpointed and
      recovered \\
    \bottomrule
  \end{tabular}
\end{table}

\begin{example}
  A client reads a \SI{12}{\kibi\byte} file in three requests of
  \SI{4}{\kibi\byte}, and the server restarts after the second
  (\cref{tab:l14-state-ex}).  The stateless server receives the file and the
  offset with every read; with a 32-byte file handle
  (\cref{sec:l14-nfs}), an 8-byte offset and a 4-byte count, each request
  carries 44 bytes of arguments.  The stateful server records the open
  file and its offset, so that a read needs only a 4-byte descriptor and
  the count, 8 bytes.  After the restart, however, the stateful server no
  longer knows descriptor~5, and the third read fails unless the state is
  recovered, whereas the stateless server serves it as if nothing had
  happened.
\end{example}

\begin{table}[htbp]
  \centering
  \caption{Reading a file from a stateless and a stateful server that
    restarts after the second read.}
  \label{tab:l14-state-ex}
  \small
  \begin{tabular}{@{}l>{\raggedright\arraybackslash}p{0.38\linewidth}>{\raggedright\arraybackslash}p{0.38\linewidth}@{}}
    \toprule
    Step & Stateless server & Stateful server \\
    \midrule
    open & \texttt{lookup(dir,"data")} returns handle $h$ &
      \texttt{open("data")} returns 5; offset 0 recorded \\
    read 1 & \texttt{read($h$,0,4096)} & \texttt{read(5,4096)}; offset
      4096 \\
    read 2 & \texttt{read($h$,4096,4096)} & \texttt{read(5,4096)}; offset
      8192 \\
    restart & nothing lost & open files and offsets lost \\
    read 3 & \texttt{read($h$,8192,4096)} served & \texttt{read(5,4096)}
      fails \\
    \bottomrule
  \end{tabular}
\end{table}

\begin{keypoint}
  Practical DFSs let clients cache blocks while interacting frequently
  with the server, a compromise between the upload/download and the remote
  access models.  The compromise requires a stateful server, which
  supports caching, locking and incremental operations at the cost of
  server resources and of recovery after failures.  A stateless server is
  simpler and survives failures by restarting, but cannot keep client
  caches consistent.
\end{keypoint}

\section{Caching state in a DFS}
\label{sec:l14-caching}

In the compromise model clients maintain a portion of the state of a file
locally, typically some of its blocks, and perform operations such as
open, read and write on this cached state.\source{P4L2, caching state in a
DFS; OSTEP ch.~49; Silberschatz et al.\ ch.~19.}  Once the same file is
cached at several places, the copies can diverge.  Suppose that clients~1
and~2 both cache file~$F$ and that client~2 modifies its copy into $F'$
(\cref{fig:l14-caching}).  The modification must reach the server, and
client~1 must learn that its copy of $F$ is stale.

\begin{figure}[htbp]
  \centering
  % Caching state in a DFS: client 2 modifies its cached copy of F, writes it
% back, and client 1 learns that its copy is stale, either by asking the
% server (client-driven) or by being notified (server-driven).
% Usage: % Caching state in a DFS: client 2 modifies its cached copy of F, writes it
% back, and client 1 learns that its copy is stale, either by asking the
% server (client-driven) or by being notified (server-driven).
% Usage: % Caching state in a DFS: client 2 modifies its cached copy of F, writes it
% back, and client 1 learns that its copy is stale, either by asking the
% server (client-driven) or by being notified (server-driven).
% Usage: \input{figures/l14-caching}   (width ~118 mm)
\begin{tikzpicture}[x=1mm, y=1mm, font=\scriptsize\sffamily,
  >={Stealth[length=1.5mm]},
  machine/.style={draw, rounded corners=0.8mm},
  doc/.style={draw, fill=white, minimum width=9mm, minimum height=6mm, inner sep=1pt},
  lab/.style={font=\scriptsize\sffamily, align=center},
  who/.style={font=\scriptsize\sffamily\bfseries, anchor=north west}]
  % client 1
  \draw[machine, fill=white] (0,25) rectangle (42,39);
  \node[who] at (1,38.5) {\iflangja{クライアント 1}{client 1}};
  \node[lab, anchor=west] at (3,30) {\iflangja{キャッシュ：}{cache:}};
  \node[doc] at (27,30) {$F$};
  \node[lab, text=ln-caution, anchor=west] at (32,30) {\iflangja{古い}{stale}};
  % client 2
  \draw[machine, fill=white] (76,25) rectangle (118,39);
  \node[who] at (77,38.5) {\iflangja{クライアント 2}{client 2}};
  \node[lab, anchor=west] at (79,30) {\iflangja{キャッシュ：}{cache:}};
  \node[doc, fill=ln-accent!20, thick] (f2) at (103,30) {$F'$};
  \node[lab, anchor=north, text=ln-accent] at (103,37.8) {\iflangja{(1) 書き込み}{(1) write}};
  % server
  \draw[machine, fill=ln-box] (38,0) rectangle (80,12);
  \node[who] at (39,11.5) {\iflangja{サーバ}{server}};
  \node[doc, fill=ln-accent!20] at (66,5) {$F'$};
  % (2) write back
  \draw[->, thick, ln-accent] (112,25) |- (80,4)
    node[lab, pos=0.25, left, text=ln-accent] {\iflangja{(2) 書き戻し}{(2) write back}};
  % (3a) client-driven validation
  \draw[->, densely dashed] (10,25) |- (38,4);
  \node[lab, anchor=south] at (24,4.6) {\iflangja{(3a) クライアント駆動：\\「$F$ は最新か？」}{(3a) client-driven:\\is $F$ current?}};
  % (3b) server-driven notification
  \draw[->, densely dashed, ln-caution] (52,12) -- (52,18) -| (34,25);
  \node[lab, anchor=west, text=ln-caution, align=left] at (53,17.5) {\iflangja{(3b) サーバ駆動：\\「$F$ は変更された」}{(3b) server-driven:\\$F$ has changed}};
\end{tikzpicture}
   (width ~118 mm)
\begin{tikzpicture}[x=1mm, y=1mm, font=\scriptsize\sffamily,
  >={Stealth[length=1.5mm]},
  machine/.style={draw, rounded corners=0.8mm},
  doc/.style={draw, fill=white, minimum width=9mm, minimum height=6mm, inner sep=1pt},
  lab/.style={font=\scriptsize\sffamily, align=center},
  who/.style={font=\scriptsize\sffamily\bfseries, anchor=north west}]
  % client 1
  \draw[machine, fill=white] (0,25) rectangle (42,39);
  \node[who] at (1,38.5) {\iflangja{クライアント 1}{client 1}};
  \node[lab, anchor=west] at (3,30) {\iflangja{キャッシュ：}{cache:}};
  \node[doc] at (27,30) {$F$};
  \node[lab, text=ln-caution, anchor=west] at (32,30) {\iflangja{古い}{stale}};
  % client 2
  \draw[machine, fill=white] (76,25) rectangle (118,39);
  \node[who] at (77,38.5) {\iflangja{クライアント 2}{client 2}};
  \node[lab, anchor=west] at (79,30) {\iflangja{キャッシュ：}{cache:}};
  \node[doc, fill=ln-accent!20, thick] (f2) at (103,30) {$F'$};
  \node[lab, anchor=north, text=ln-accent] at (103,37.8) {\iflangja{(1) 書き込み}{(1) write}};
  % server
  \draw[machine, fill=ln-box] (38,0) rectangle (80,12);
  \node[who] at (39,11.5) {\iflangja{サーバ}{server}};
  \node[doc, fill=ln-accent!20] at (66,5) {$F'$};
  % (2) write back
  \draw[->, thick, ln-accent] (112,25) |- (80,4)
    node[lab, pos=0.25, left, text=ln-accent] {\iflangja{(2) 書き戻し}{(2) write back}};
  % (3a) client-driven validation
  \draw[->, densely dashed] (10,25) |- (38,4);
  \node[lab, anchor=south] at (24,4.6) {\iflangja{(3a) クライアント駆動：\\「$F$ は最新か？」}{(3a) client-driven:\\is $F$ current?}};
  % (3b) server-driven notification
  \draw[->, densely dashed, ln-caution] (52,12) -- (52,18) -| (34,25);
  \node[lab, anchor=west, text=ln-caution, align=left] at (53,17.5) {\iflangja{(3b) サーバ駆動：\\「$F$ は変更された」}{(3b) server-driven:\\$F$ has changed}};
\end{tikzpicture}
   (width ~118 mm)
\begin{tikzpicture}[x=1mm, y=1mm, font=\scriptsize\sffamily,
  >={Stealth[length=1.5mm]},
  machine/.style={draw, rounded corners=0.8mm},
  doc/.style={draw, fill=white, minimum width=9mm, minimum height=6mm, inner sep=1pt},
  lab/.style={font=\scriptsize\sffamily, align=center},
  who/.style={font=\scriptsize\sffamily\bfseries, anchor=north west}]
  % client 1
  \draw[machine, fill=white] (0,25) rectangle (42,39);
  \node[who] at (1,38.5) {\iflangja{クライアント 1}{client 1}};
  \node[lab, anchor=west] at (3,30) {\iflangja{キャッシュ：}{cache:}};
  \node[doc] at (27,30) {$F$};
  \node[lab, text=ln-caution, anchor=west] at (32,30) {\iflangja{古い}{stale}};
  % client 2
  \draw[machine, fill=white] (76,25) rectangle (118,39);
  \node[who] at (77,38.5) {\iflangja{クライアント 2}{client 2}};
  \node[lab, anchor=west] at (79,30) {\iflangja{キャッシュ：}{cache:}};
  \node[doc, fill=ln-accent!20, thick] (f2) at (103,30) {$F'$};
  \node[lab, anchor=north, text=ln-accent] at (103,37.8) {\iflangja{(1) 書き込み}{(1) write}};
  % server
  \draw[machine, fill=ln-box] (38,0) rectangle (80,12);
  \node[who] at (39,11.5) {\iflangja{サーバ}{server}};
  \node[doc, fill=ln-accent!20] at (66,5) {$F'$};
  % (2) write back
  \draw[->, thick, ln-accent] (112,25) |- (80,4)
    node[lab, pos=0.25, left, text=ln-accent] {\iflangja{(2) 書き戻し}{(2) write back}};
  % (3a) client-driven validation
  \draw[->, densely dashed] (10,25) |- (38,4);
  \node[lab, anchor=south] at (24,4.6) {\iflangja{(3a) クライアント駆動：\\「$F$ は最新か？」}{(3a) client-driven:\\is $F$ current?}};
  % (3b) server-driven notification
  \draw[->, densely dashed, ln-caution] (52,12) -- (52,18) -| (34,25);
  \node[lab, anchor=west, text=ln-caution, align=left] at (53,17.5) {\iflangja{(3b) サーバ駆動：\\「$F$ は変更された」}{(3b) server-driven:\\$F$ has changed}};
\end{tikzpicture}

  \caption{Client 2 modifies its cached copy of $F$ and writes it back.
    Client 1 learns that its copy is stale either by asking the server
    (client-driven) or by being notified (server-driven).}
  \label{fig:l14-caching}
\end{figure}

The problem is that of cache coherence\index{cache coherence} on shared
memory multiprocessors (Lesson~10), where a write by one CPU makes the
copies of a memory location in the caches of the other CPUs stale and the
hardware updates or invalidates them.  In a DFS the coherence mechanism is
implemented in software, and it is either
\begin{description}
  \item[client-driven] the client asks the server, when it opens a file or
    periodically, whether its cached data are still current; or
  \item[server-driven] the server, which knows which clients cache a file,
    notifies them when the file changes.
\end{description}
Client-driven checks keep the server simple but cost a request each and
leave a client with stale data until its next check; server-driven
notification requires a stateful server.  When and how often either
mechanism acts is determined by the file sharing semantics that the DFS
provides (\cref{sec:l14-semantics}).

In a DFS with one server and many clients, files or their blocks can be
cached in three places:
\begin{itemize}
  \item in the memory of a client, like the buffer cache of Lesson~11: the
    fastest location, but limited in size and lost when the client
    crashes;
  \item on a storage device of the client, a hard disk or SSD: slower than
    memory but faster than the network, larger, and persistent;
  \item in the buffer cache in the memory of the server, which saves the
    server disk accesses.  How much it helps depends on the load: when many
    clients access different files and their requests interleave, the data
    of one client displace those of others and the hit rate falls.
\end{itemize}

\section{File sharing semantics}
\label{sec:l14-semantics}

On a single machine, when process~A writes to a file that process~B also
has open, B's next read returns the new data: both processes access the
file through the same buffer cache, so every write is visible
immediately.\source{P4L2, file sharing semantics on a DFS; Silberschatz et
al.\ ch.~15; OSTEP ch.~49.}  In a DFS whose clients cache data, a write by
client~A first changes only A's cache.  It reaches the server when A sends
it there, and client~B sees it only after B has obtained it from the
server.  The guarantees about when one client observes the writes of
another form the \emph{file sharing semantics} of a DFS.

\begin{definition}[UNIX semantics]
  Under \term{UNIX semantics}[UNIXセマンティクス] every write to a file is
  visible immediately to every subsequent read, on any machine.
\end{definition}

UNIX semantics is what applications expect from a local file system.  In
a DFS it requires every write to be propagated before any other client
reads the file, which makes it expensive to provide.

\begin{definition}[Session semantics]
  Under \term{session semantics}[セッションセマンティクス] a client obtains
  the current version of a file from the server when it opens the file,
  and writes its changes back to the server when it closes the file.  The
  period between an open and the corresponding close is a \emph{session};
  changes made by other clients during a session are not visible in it.
\end{definition}

Session semantics is easy to reason about, but it may be insufficient.
Clients that share a file while both have it open do not see each other's
changes, and the longer the sessions last, the longer the copies of
different clients remain inconsistent.

With \emph{periodic updates} a client writes its changes back to the
server periodically, and cached data are invalidated or revalidated
periodically.  The period bounds the time during which clients can observe
inconsistent data, and the DFS may offer explicit operations such as
\texttt{flush} or \texttt{sync} with which an application forces its changes
to the server at once.  A client's right to use its cached data is thus
limited in time: the client holds a lease on the data.

\begin{definition}[Lease]
  A \term{lease}[リース] is a right that a server grants to a client for a
  limited period of time, such as the right to use cached data or to hold
  a lock.  When the period expires, the right is lost unless the client has
  renewed it.
\end{definition}

A lease on cached data need not be exclusive: several clients can hold
leases on the same file at the same time.  Two further approaches avoid
inconsistent views altogether.

\begin{definition}[Immutable file]
  An \term{immutable file}[不変ファイル] is never modified after it has
  been created.  To change its contents, a new file is created, and the old
  one may be deleted.
\end{definition}

All copies of an immutable file are identical, so sharing it cannot
produce inconsistency.  Services that store photos, for example, typically
do not modify an uploaded photo in place; an edited photo is stored as a
new file.  Finally, a DFS can offer \emph{transactions}: an application
groups several updates, of one or several files, into a transaction, and
the DFS applies them atomically---either all take effect or none
does---as a database system does.

\begin{example}
  \label{ex:l14-semantics}
  File $F$ initially contains \texttt{x}.  Client~A opens $F$ at
  $t=\SI{0}{\second}$, writes \texttt{y} at \SI{4}{\second} and closes $F$
  at \SI{10}{\second}.  Client~B opens $F$ at \SI{2}{\second}, reads it at
  \SI{6}{\second} and \SI{12}{\second} and closes it at \SI{14}{\second};
  it opens $F$ again at \SI{15}{\second} and reads it at \SI{16}{\second}
  (\cref{fig:l14-semantics}).  For periodic updates, assume that A's client
  writes back dirty data every \SI{5}{\second}, at 5, 10, 15,
  \ldots~seconds, and that B's client revalidates its cache at 3, 8, 13,
  \ldots~seconds.  \Cref{tab:l14-semantics} lists what B reads.  Under
  session semantics B sees \texttt{y} only in its second session.  Under
  periodic updates it sees \texttt{y} after its first revalidation that
  follows A's write-back, \SI{4}{\second} after the write; in general no
  write stays invisible to B for longer than the sum of the two periods,
  \SI{10}{\second}.
\end{example}

\begin{figure}[htbp]
  \centering
  % File sharing semantics: client A writes y into F during a session, client
% B reads F.  (a) Session semantics: write-back on close, fetch on open.
% (b) Periodic updates: A writes back every 5 s, B revalidates at 3, 8, 13 s.
% Usage: % File sharing semantics: client A writes y into F during a session, client
% B reads F.  (a) Session semantics: write-back on close, fetch on open.
% (b) Periodic updates: A writes back every 5 s, B revalidates at 3, 8, 13 s.
% Usage: % File sharing semantics: client A writes y into F during a session, client
% B reads F.  (a) Session semantics: write-back on close, fetch on open.
% (b) Periodic updates: A writes back every 5 s, B revalidates at 3, 8, 13 s.
% Usage: \input{figures/l14-semantics}   (width ~116 mm)
\begin{tikzpicture}[x=1mm, y=1mm, font=\scriptsize\sffamily,
  >={Stealth[length=1.3mm]},
  row/.style={font=\scriptsize\sffamily, anchor=east},
  lab/.style={font=\scriptsize\sffamily, inner sep=0.8pt},
  val/.style={font=\scriptsize\ttfamily, inner sep=0.8pt},
  sess/.style={draw, fill=ln-box},
  ev/.style={circle, fill=black, inner sep=0pt, minimum size=1.3mm},
  ttl/.style={font=\footnotesize\sffamily\bfseries, anchor=west}]
  % time t (s) -> x = 20 + 5.6 t (mm)
  % \lpanel{title}  rows: A at 14, server at 7, B at 0 (inside a scope)
  \def\lpanel#1{%
    \node[ttl] at (0,21) {#1};
    \node[row] at (18,14) {\iflangja{クライアント A}{client A}};
    \node[row] at (18,7) {\iflangja{サーバ}{server}};
    \node[row] at (18,0) {\iflangja{クライアント B}{client B}};
    \foreach \y in {0,7,14} { \draw[ln-rule] (20,\y) -- (116,\y); }
    % sessions
    \draw[sess] ({20+5.6*0},12.6) rectangle ({20+5.6*10},15.4);
    \draw[sess] ({20+5.6*2},-1.4) rectangle ({20+5.6*14},1.4);
    \draw[sess] ({20+5.6*15},-1.4) rectangle ({20+5.6*17},1.4);
    % A writes y at 4 s
    \node[ev] at ({20+5.6*4},14) {};
    \node[lab, anchor=south] at ({20+5.6*4},15.6) {\iflangja{\texttt{y} を書く}{write \texttt{y}}};
    % B reads at 6, 12, 16 s
    \foreach \t in {6,12,16} { \node[ev] at ({20+5.6*\t},0) {}; }
  }
  % ---------------- (a) session semantics
  \begin{scope}[shift={(0,33)}]
    \lpanel{\iflangja{(a) セッションセマンティクス}{(a) session semantics}}
    \draw[->, thick, ln-accent] ({20+5.6*10},12.6) -- node[lab, right, text=ln-accent] {\iflangja{close で書き戻し}{write back on close}} ({20+5.6*10},7.6);
    \draw[->, thick, ln-accent] ({20+5.6*2},6.4) -- node[val, right] {x} ({20+5.6*2},1.6);
    \draw[->, thick, ln-accent] ({20+5.6*15},6.4) -- node[val, left] {y} ({20+5.6*15},1.6);
    \foreach \t/\v in {6/x,12/x,16/y} { \node[lab, anchor=north] at ({20+5.6*\t},-1.8) {\iflangja{読む：}{read }\texttt{\v}}; }
  \end{scope}
  % ---------------- (b) periodic updates
  \begin{scope}[shift={(0,0)}]
    \lpanel{\iflangja{(b) 周期的な更新（5 秒）}{(b) periodic updates (5 s)}}
    \draw[->, thick, ln-accent] ({20+5.6*5},12.6) -- node[lab, right, text=ln-accent] {\iflangja{書き戻し}{write back}} ({20+5.6*5},7.6);
    \foreach \t/\v in {3/x,8/y,13/y} {
      \draw[->, densely dashed] ({20+5.6*\t},6.4) -- node[val, right] {\v} ({20+5.6*\t},1.6);
    }
    \foreach \t/\v in {6/x,12/y,16/y} { \node[lab, anchor=north] at ({20+5.6*\t},-1.8) {\iflangja{読む：}{read }\texttt{\v}}; }
  \end{scope}
  % time axis
  \draw[->] (20,-9) -- (117,-9);
  \foreach \t in {0,5,10,15} {
    \draw ({20+5.6*\t},-8.2) -- ({20+5.6*\t},-9.8);
    \node[lab, anchor=north] at ({20+5.6*\t},-10) {\t};
  }
  \node[lab, anchor=south east] at (117,-8.6) {\iflangja{時間（秒）}{time (s)}};
\end{tikzpicture}
   (width ~116 mm)
\begin{tikzpicture}[x=1mm, y=1mm, font=\scriptsize\sffamily,
  >={Stealth[length=1.3mm]},
  row/.style={font=\scriptsize\sffamily, anchor=east},
  lab/.style={font=\scriptsize\sffamily, inner sep=0.8pt},
  val/.style={font=\scriptsize\ttfamily, inner sep=0.8pt},
  sess/.style={draw, fill=ln-box},
  ev/.style={circle, fill=black, inner sep=0pt, minimum size=1.3mm},
  ttl/.style={font=\footnotesize\sffamily\bfseries, anchor=west}]
  % time t (s) -> x = 20 + 5.6 t (mm)
  % \lpanel{title}  rows: A at 14, server at 7, B at 0 (inside a scope)
  \def\lpanel#1{%
    \node[ttl] at (0,21) {#1};
    \node[row] at (18,14) {\iflangja{クライアント A}{client A}};
    \node[row] at (18,7) {\iflangja{サーバ}{server}};
    \node[row] at (18,0) {\iflangja{クライアント B}{client B}};
    \foreach \y in {0,7,14} { \draw[ln-rule] (20,\y) -- (116,\y); }
    % sessions
    \draw[sess] ({20+5.6*0},12.6) rectangle ({20+5.6*10},15.4);
    \draw[sess] ({20+5.6*2},-1.4) rectangle ({20+5.6*14},1.4);
    \draw[sess] ({20+5.6*15},-1.4) rectangle ({20+5.6*17},1.4);
    % A writes y at 4 s
    \node[ev] at ({20+5.6*4},14) {};
    \node[lab, anchor=south] at ({20+5.6*4},15.6) {\iflangja{\texttt{y} を書く}{write \texttt{y}}};
    % B reads at 6, 12, 16 s
    \foreach \t in {6,12,16} { \node[ev] at ({20+5.6*\t},0) {}; }
  }
  % ---------------- (a) session semantics
  \begin{scope}[shift={(0,33)}]
    \lpanel{\iflangja{(a) セッションセマンティクス}{(a) session semantics}}
    \draw[->, thick, ln-accent] ({20+5.6*10},12.6) -- node[lab, right, text=ln-accent] {\iflangja{close で書き戻し}{write back on close}} ({20+5.6*10},7.6);
    \draw[->, thick, ln-accent] ({20+5.6*2},6.4) -- node[val, right] {x} ({20+5.6*2},1.6);
    \draw[->, thick, ln-accent] ({20+5.6*15},6.4) -- node[val, left] {y} ({20+5.6*15},1.6);
    \foreach \t/\v in {6/x,12/x,16/y} { \node[lab, anchor=north] at ({20+5.6*\t},-1.8) {\iflangja{読む：}{read }\texttt{\v}}; }
  \end{scope}
  % ---------------- (b) periodic updates
  \begin{scope}[shift={(0,0)}]
    \lpanel{\iflangja{(b) 周期的な更新（5 秒）}{(b) periodic updates (5 s)}}
    \draw[->, thick, ln-accent] ({20+5.6*5},12.6) -- node[lab, right, text=ln-accent] {\iflangja{書き戻し}{write back}} ({20+5.6*5},7.6);
    \foreach \t/\v in {3/x,8/y,13/y} {
      \draw[->, densely dashed] ({20+5.6*\t},6.4) -- node[val, right] {\v} ({20+5.6*\t},1.6);
    }
    \foreach \t/\v in {6/x,12/y,16/y} { \node[lab, anchor=north] at ({20+5.6*\t},-1.8) {\iflangja{読む：}{read }\texttt{\v}}; }
  \end{scope}
  % time axis
  \draw[->] (20,-9) -- (117,-9);
  \foreach \t in {0,5,10,15} {
    \draw ({20+5.6*\t},-8.2) -- ({20+5.6*\t},-9.8);
    \node[lab, anchor=north] at ({20+5.6*\t},-10) {\t};
  }
  \node[lab, anchor=south east] at (117,-8.6) {\iflangja{時間（秒）}{time (s)}};
\end{tikzpicture}
   (width ~116 mm)
\begin{tikzpicture}[x=1mm, y=1mm, font=\scriptsize\sffamily,
  >={Stealth[length=1.3mm]},
  row/.style={font=\scriptsize\sffamily, anchor=east},
  lab/.style={font=\scriptsize\sffamily, inner sep=0.8pt},
  val/.style={font=\scriptsize\ttfamily, inner sep=0.8pt},
  sess/.style={draw, fill=ln-box},
  ev/.style={circle, fill=black, inner sep=0pt, minimum size=1.3mm},
  ttl/.style={font=\footnotesize\sffamily\bfseries, anchor=west}]
  % time t (s) -> x = 20 + 5.6 t (mm)
  % \lpanel{title}  rows: A at 14, server at 7, B at 0 (inside a scope)
  \def\lpanel#1{%
    \node[ttl] at (0,21) {#1};
    \node[row] at (18,14) {\iflangja{クライアント A}{client A}};
    \node[row] at (18,7) {\iflangja{サーバ}{server}};
    \node[row] at (18,0) {\iflangja{クライアント B}{client B}};
    \foreach \y in {0,7,14} { \draw[ln-rule] (20,\y) -- (116,\y); }
    % sessions
    \draw[sess] ({20+5.6*0},12.6) rectangle ({20+5.6*10},15.4);
    \draw[sess] ({20+5.6*2},-1.4) rectangle ({20+5.6*14},1.4);
    \draw[sess] ({20+5.6*15},-1.4) rectangle ({20+5.6*17},1.4);
    % A writes y at 4 s
    \node[ev] at ({20+5.6*4},14) {};
    \node[lab, anchor=south] at ({20+5.6*4},15.6) {\iflangja{\texttt{y} を書く}{write \texttt{y}}};
    % B reads at 6, 12, 16 s
    \foreach \t in {6,12,16} { \node[ev] at ({20+5.6*\t},0) {}; }
  }
  % ---------------- (a) session semantics
  \begin{scope}[shift={(0,33)}]
    \lpanel{\iflangja{(a) セッションセマンティクス}{(a) session semantics}}
    \draw[->, thick, ln-accent] ({20+5.6*10},12.6) -- node[lab, right, text=ln-accent] {\iflangja{close で書き戻し}{write back on close}} ({20+5.6*10},7.6);
    \draw[->, thick, ln-accent] ({20+5.6*2},6.4) -- node[val, right] {x} ({20+5.6*2},1.6);
    \draw[->, thick, ln-accent] ({20+5.6*15},6.4) -- node[val, left] {y} ({20+5.6*15},1.6);
    \foreach \t/\v in {6/x,12/x,16/y} { \node[lab, anchor=north] at ({20+5.6*\t},-1.8) {\iflangja{読む：}{read }\texttt{\v}}; }
  \end{scope}
  % ---------------- (b) periodic updates
  \begin{scope}[shift={(0,0)}]
    \lpanel{\iflangja{(b) 周期的な更新（5 秒）}{(b) periodic updates (5 s)}}
    \draw[->, thick, ln-accent] ({20+5.6*5},12.6) -- node[lab, right, text=ln-accent] {\iflangja{書き戻し}{write back}} ({20+5.6*5},7.6);
    \foreach \t/\v in {3/x,8/y,13/y} {
      \draw[->, densely dashed] ({20+5.6*\t},6.4) -- node[val, right] {\v} ({20+5.6*\t},1.6);
    }
    \foreach \t/\v in {6/x,12/y,16/y} { \node[lab, anchor=north] at ({20+5.6*\t},-1.8) {\iflangja{読む：}{read }\texttt{\v}}; }
  \end{scope}
  % time axis
  \draw[->] (20,-9) -- (117,-9);
  \foreach \t in {0,5,10,15} {
    \draw ({20+5.6*\t},-8.2) -- ({20+5.6*\t},-9.8);
    \node[lab, anchor=north] at ({20+5.6*\t},-10) {\t};
  }
  \node[lab, anchor=south east] at (117,-8.6) {\iflangja{時間（秒）}{time (s)}};
\end{tikzpicture}

  \caption{The scenario of \cref{ex:l14-semantics}.  Shaded bars are
    sessions.  (a) Under session semantics A writes back at close and B
    fetches the file at open.  (b) Under periodic updates A writes back at
    \SI{5}{\second} and B revalidates every \SI{5}{\second} (dashed).}
  \label{fig:l14-semantics}
\end{figure}

\begin{table}[htbp]
  \centering
  \caption{Contents of $F$ read by client B in \cref{ex:l14-semantics}.}
  \label{tab:l14-semantics}
  \small
  \begin{tabular}{@{}lccc@{}}
    \toprule
    Read by B at & UNIX semantics & Session semantics & Periodic updates \\
    \midrule
    \SI{6}{\second}  & \texttt{y} & \texttt{x} & \texttt{x} \\
    \SI{12}{\second} & \texttt{y} & \texttt{x} & \texttt{y} \\
    \SI{16}{\second} & \texttt{y} & \texttt{y} & \texttt{y} \\
    \bottomrule
  \end{tabular}
\end{table}

\needspace{6\baselineskip}
\section{Files and directories}
\label{sec:l14-dirs}

No single semantics is the best choice for every DFS.\source{P4L2, file
vs.\ directory service; Silberschatz et al.\ ch.~19.}  The right choice
depends on the workload: how often files are shared, how often they are
written, and how important a consistent view is to the applications.
Knowing these access patterns allows the designer to optimize for the
common case\index{optimize for the common case} (Lesson~1).

A DFS serves two kinds of files whose workloads differ.  A regular file is
typically used by one client at a time, and a delay before other clients
see its changes is rarely noticed.  A directory is shared far more
often---every client that creates, removes or renames a file in a shared
directory modifies it---and an inconsistent view of a directory is
harmful: two clients could, for instance, both create a file of the same
name.  A DFS can therefore apply different policies to the two, for
example session semantics for files and UNIX semantics for directories,
or periodic updates with less frequent write-back for files than for
directories.

\begin{keypoint}
  File sharing semantics define when one client observes the writes of
  another.  UNIX semantics makes every write visible immediately but is
  expensive in a DFS; session semantics propagates changes at close and
  open; periodic updates bound the inconsistency by the update period;
  immutable files and transactions avoid inconsistent views altogether.
  The choice follows the workload, and files and directories may be
  treated differently.
\end{keypoint}

\section{Replication and partitioning}
\label{sec:l14-replication}

A file service that runs on several machines distributes the files among
them in one of two ways (\cref{fig:l14-models}b and~c).\source{P4L2,
replication and partitioning; Silberschatz et al.\ ch.~19.}

\begin{definition}[Replication]
  With \term{replication}[レプリケーション] every server machine holds a
  complete copy, a \emph{replica}, of all files.
\end{definition}

Replication balances the load, since a read can be served by any replica,
and it provides availability and fault tolerance: when a machine fails,
the others continue to serve all files.  Writes, however, become more
complex.  Either every write is performed synchronously on all replicas,
which makes writes slow, or it is performed on one replica and propagated
to the others afterwards, in which case the replicas differ for a while.
Replicas that have diverged must be reconciled, for example by voting, in
which the contents held by a majority of the replicas prevail.

\begin{definition}[Partitioning]
  With \term{partitioning}[パーティショニング] each server machine holds a
  subset of the files, and the subsets together make up the file system.
\end{definition}

Partitioning improves availability compared with a single server, since
the failure of a machine affects only the files it holds.  It scales with
the size of the file system: more files need only more machines.  A write
to a single file involves only the machine that holds it, so writes remain
simple.  On the other hand, the portion of the data held by a failed
machine is lost, and balancing the load is harder: if popular files are
concentrated on a few machines, these become hot spots.
\Cref{tab:l14-repl} compares the two approaches.  They can also be
combined: the files are partitioned, and every partition is replicated on
several machines.%
\begin{marginfigure}
  \centering
  % Replication combined with partitioning: two partitions of 250 files, each
% replicated on two of four machines.
% Usage: % Replication combined with partitioning: two partitions of 250 files, each
% replicated on two of four machines.
% Usage: % Replication combined with partitioning: two partitions of 250 files, each
% replicated on two of four machines.
% Usage: \input{figures/l14-combined}   (margin figure, width ~48 mm)
\begin{tikzpicture}[x=1mm, y=1mm, font=\scriptsize\sffamily,
  sv/.style={draw, fill=ln-box, minimum width=18mm, minimum height=8mm, inner sep=1pt, align=center},
  grp/.style={draw, densely dashed, ln-muted, rounded corners=0.8mm},
  lab/.style={font=\scriptsize\sffamily, text=ln-muted, anchor=south west, inner sep=0.5pt}]
  \foreach \y/\p/\f/\a/\b in {16/1/1--250/1/2, 0/2/251--500/3/4} {
    \draw[grp] (0,\y-5) rectangle (47,\y+9);
    \node[lab] at (1,\y+5) {\iflangja{パーティション}{partition} \p: \iflangja{ファイル}{files} \f};
    \node[sv] at (11.5,\y) {S\a};
    \node[sv] at (35.5,\y) {S\b};
  }
\end{tikzpicture}
   (margin figure, width ~48 mm)
\begin{tikzpicture}[x=1mm, y=1mm, font=\scriptsize\sffamily,
  sv/.style={draw, fill=ln-box, minimum width=18mm, minimum height=8mm, inner sep=1pt, align=center},
  grp/.style={draw, densely dashed, ln-muted, rounded corners=0.8mm},
  lab/.style={font=\scriptsize\sffamily, text=ln-muted, anchor=south west, inner sep=0.5pt}]
  \foreach \y/\p/\f/\a/\b in {16/1/1--250/1/2, 0/2/251--500/3/4} {
    \draw[grp] (0,\y-5) rectangle (47,\y+9);
    \node[lab] at (1,\y+5) {\iflangja{パーティション}{partition} \p: \iflangja{ファイル}{files} \f};
    \node[sv] at (11.5,\y) {S\a};
    \node[sv] at (35.5,\y) {S\b};
  }
\end{tikzpicture}
   (margin figure, width ~48 mm)
\begin{tikzpicture}[x=1mm, y=1mm, font=\scriptsize\sffamily,
  sv/.style={draw, fill=ln-box, minimum width=18mm, minimum height=8mm, inner sep=1pt, align=center},
  grp/.style={draw, densely dashed, ln-muted, rounded corners=0.8mm},
  lab/.style={font=\scriptsize\sffamily, text=ln-muted, anchor=south west, inner sep=0.5pt}]
  \foreach \y/\p/\f/\a/\b in {16/1/1--250/1/2, 0/2/251--500/3/4} {
    \draw[grp] (0,\y-5) rectangle (47,\y+9);
    \node[lab] at (1,\y+5) {\iflangja{パーティション}{partition} \p: \iflangja{ファイル}{files} \f};
    \node[sv] at (11.5,\y) {S\a};
    \node[sv] at (35.5,\y) {S\b};
  }
\end{tikzpicture}

  \caption{Two partitions, each replicated on two machines.}
  \label{fig:l14-combined}
\end{marginfigure}

\begin{table}[htbp]
  \centering
  \caption{Replication and partitioning.}
  \label{tab:l14-repl}
  \small
  \begin{tabular}{@{}>{\raggedright\arraybackslash}p{0.2\linewidth}>{\raggedright\arraybackslash}p{0.36\linewidth}>{\raggedright\arraybackslash}p{0.36\linewidth}@{}}
    \toprule
    & Replication & Partitioning \\
    \midrule
    Capacity & that of one machine & sum over all machines \\
    Reads & served by any replica; load balanced & served by the machine
      holding the file; hot spots possible \\
    Writes & on all replicas, synchronously or propagated later; replicas
      must be reconciled & on one machine; simple \\
    Machine failure & no files lost & the files of that machine lost \\
    \bottomrule
  \end{tabular}
\end{table}

\begin{example}
  Four server machines can store 250 files each.  Replicated, the DFS
  holds 250 files, and the failure of one machine---indeed of any
  three---loses none of them.  Partitioned, it holds \num{1000} files, and
  the failure of one machine loses $250/1000 = \SI{25}{\percent}$ of them.
  Combined as in \cref{fig:l14-combined}, with two partitions of 250 files
  each replicated on two machines, the DFS holds 500 files; the failure of
  one machine loses none of them, and only the failure of both machines of
  the same partition loses \SI{50}{\percent}.
\end{example}

\section{The Network File System}
\label{sec:l14-nfs}

The \term{Network File System}[NFS]\index{NFS|see{Network File System}}
(NFS), developed at Sun Microsystems in the 1980s \cite{sandberg1985}, is a
widely used DFS in which clients access files that reside on a remote
server.\source{P4L2, NFS design; Sandberg et al.\ \cite{sandberg1985};
OSTEP ch.~49; Silberschatz et al.\ ch.~15.}  Its structure is shown in
\cref{fig:l14-nfs}.  An application on the client performs ordinary file
system calls.  The VFS finds that the file belongs to an NFS file system
and passes the operation to the \emph{NFS client}, a file system
implementation that turns the operation into a remote procedure call to
the server.  There the \emph{NFS server} component receives the request,
performs the operation on the server's files through the server's own VFS
and local file system, and returns the result, which travels back to the
application.  NFS is built on Sun RPC\index{Sun RPC} and its data
representation XDR (Lesson~13).

\begin{figure}[htbp]
  \centering
  % NFS architecture: the VFS on the client passes operations on NFS files to
% the NFS client, which sends RPC requests to the NFS server; the server
% performs them through its own VFS and local file system.
% Usage: % NFS architecture: the VFS on the client passes operations on NFS files to
% the NFS client, which sends RPC requests to the NFS server; the server
% performs them through its own VFS and local file system.
% Usage: % NFS architecture: the VFS on the client passes operations on NFS files to
% the NFS client, which sends RPC requests to the NFS server; the server
% performs them through its own VFS and local file system.
% Usage: \input{figures/l14-nfs}   (width ~118 mm)
\begin{tikzpicture}[x=1mm, y=1mm, font=\scriptsize\sffamily,
  >={Stealth[length=1.5mm]},
  machine/.style={draw, rounded corners=0.8mm, fill=white},
  box/.style={draw, fill=white, minimum height=6mm, inner sep=1pt, align=center},
  kbox/.style={draw, fill=ln-box, minimum height=6mm, inner sep=1pt, align=center},
  nfs/.style={draw, fill=ln-accent!18, minimum height=6mm, inner sep=1pt, align=center},
  dev/.style={draw, fill=ln-accent!35, minimum height=5mm, inner sep=1pt, align=center},
  lab/.style={font=\scriptsize\sffamily, align=center},
  who/.style={font=\footnotesize\sffamily\bfseries, anchor=north west}]
  % ---------------- client machine
  \draw[machine] (0,0) rectangle (56,68);
  \node[who] at (1,67.5) {\iflangja{クライアントマシン}{client machine}};
  \node[box, minimum width=44mm] (app) at (28,56) {\iflangja{アプリケーション}{applications}};
  \draw[densely dashed, ln-muted] (1,50.5) -- (55,50.5);
  \node[kbox, minimum width=44mm] (sc) at (28,45) {\iflangja{システムコール層}{system call layer}};
  \node[kbox, minimum width=44mm] (vfs) at (28,35) {VFS};
  \node[kbox, minimum width=20mm] (lfs) at (16,24) {\iflangja{ローカル FS}{local FS}};
  \node[nfs, minimum width=20mm] (nc) at (40,24) {\iflangja{NFS クライアント}{NFS client}};
  \node[dev, minimum width=20mm] (d1) at (16,12) {\iflangja{ディスク}{disk}};
  \node[nfs, minimum width=20mm] (cs) at (40,12) {\iflangja{RPC スタブ}{RPC stub}};
  \draw[<->, thick] (app) -- (sc);
  \draw[<->, thick] (sc) -- (vfs);
  \draw[<->, thick] (lfs.north |- vfs.south) -- (lfs.north);
  \draw[<->, thick] (nc.north |- vfs.south) -- (nc.north);
  \draw[<->, thick] (lfs) -- (d1);
  \draw[<->, thick] (nc) -- (cs);
  % ---------------- server machine
  \draw[machine] (62,0) rectangle (118,68);
  \node[who] at (63,67.5) {\iflangja{サーバマシン}{server machine}};
  \node[nfs, minimum width=44mm] (ss) at (90,56) {\iflangja{RPC スタブ}{RPC stub}};
  \node[nfs, minimum width=44mm] (ns) at (90,45) {\iflangja{NFS サーバ}{NFS server}};
  \node[kbox, minimum width=44mm] (svfs) at (90,35) {VFS};
  \node[kbox, minimum width=44mm] (slfs) at (90,24) {\iflangja{ローカル FS}{local FS}};
  \node[dev, minimum width=44mm] (d2) at (90,12) {\iflangja{ディスク}{disk}};
  \draw[<->, thick] (ss) -- (ns);
  \draw[<->, thick] (ns) -- (svfs);
  \draw[<->, thick] (svfs) -- (slfs);
  \draw[<->, thick] (slfs) -- (d2);
  % ---------------- network
  \draw[<->, line width=1.4pt, ln-accent] (cs.east) -- (59,12) |- (ss.west);
  \node[lab, rotate=90, fill=white, inner sep=0.8pt, text=ln-accent] at (59,34) {\iflangja{ネットワーク（RPC）}{network (RPC)}};
  \node[lab, text=ln-muted, align=left, anchor=west] at (66,4.2)
    {\iflangja{要求はファイルハンドルで対象ファイルを指定}{requests name files by their file handles}};
\end{tikzpicture}
   (width ~118 mm)
\begin{tikzpicture}[x=1mm, y=1mm, font=\scriptsize\sffamily,
  >={Stealth[length=1.5mm]},
  machine/.style={draw, rounded corners=0.8mm, fill=white},
  box/.style={draw, fill=white, minimum height=6mm, inner sep=1pt, align=center},
  kbox/.style={draw, fill=ln-box, minimum height=6mm, inner sep=1pt, align=center},
  nfs/.style={draw, fill=ln-accent!18, minimum height=6mm, inner sep=1pt, align=center},
  dev/.style={draw, fill=ln-accent!35, minimum height=5mm, inner sep=1pt, align=center},
  lab/.style={font=\scriptsize\sffamily, align=center},
  who/.style={font=\footnotesize\sffamily\bfseries, anchor=north west}]
  % ---------------- client machine
  \draw[machine] (0,0) rectangle (56,68);
  \node[who] at (1,67.5) {\iflangja{クライアントマシン}{client machine}};
  \node[box, minimum width=44mm] (app) at (28,56) {\iflangja{アプリケーション}{applications}};
  \draw[densely dashed, ln-muted] (1,50.5) -- (55,50.5);
  \node[kbox, minimum width=44mm] (sc) at (28,45) {\iflangja{システムコール層}{system call layer}};
  \node[kbox, minimum width=44mm] (vfs) at (28,35) {VFS};
  \node[kbox, minimum width=20mm] (lfs) at (16,24) {\iflangja{ローカル FS}{local FS}};
  \node[nfs, minimum width=20mm] (nc) at (40,24) {\iflangja{NFS クライアント}{NFS client}};
  \node[dev, minimum width=20mm] (d1) at (16,12) {\iflangja{ディスク}{disk}};
  \node[nfs, minimum width=20mm] (cs) at (40,12) {\iflangja{RPC スタブ}{RPC stub}};
  \draw[<->, thick] (app) -- (sc);
  \draw[<->, thick] (sc) -- (vfs);
  \draw[<->, thick] (lfs.north |- vfs.south) -- (lfs.north);
  \draw[<->, thick] (nc.north |- vfs.south) -- (nc.north);
  \draw[<->, thick] (lfs) -- (d1);
  \draw[<->, thick] (nc) -- (cs);
  % ---------------- server machine
  \draw[machine] (62,0) rectangle (118,68);
  \node[who] at (63,67.5) {\iflangja{サーバマシン}{server machine}};
  \node[nfs, minimum width=44mm] (ss) at (90,56) {\iflangja{RPC スタブ}{RPC stub}};
  \node[nfs, minimum width=44mm] (ns) at (90,45) {\iflangja{NFS サーバ}{NFS server}};
  \node[kbox, minimum width=44mm] (svfs) at (90,35) {VFS};
  \node[kbox, minimum width=44mm] (slfs) at (90,24) {\iflangja{ローカル FS}{local FS}};
  \node[dev, minimum width=44mm] (d2) at (90,12) {\iflangja{ディスク}{disk}};
  \draw[<->, thick] (ss) -- (ns);
  \draw[<->, thick] (ns) -- (svfs);
  \draw[<->, thick] (svfs) -- (slfs);
  \draw[<->, thick] (slfs) -- (d2);
  % ---------------- network
  \draw[<->, line width=1.4pt, ln-accent] (cs.east) -- (59,12) |- (ss.west);
  \node[lab, rotate=90, fill=white, inner sep=0.8pt, text=ln-accent] at (59,34) {\iflangja{ネットワーク（RPC）}{network (RPC)}};
  \node[lab, text=ln-muted, align=left, anchor=west] at (66,4.2)
    {\iflangja{要求はファイルハンドルで対象ファイルを指定}{requests name files by their file handles}};
\end{tikzpicture}
   (width ~118 mm)
\begin{tikzpicture}[x=1mm, y=1mm, font=\scriptsize\sffamily,
  >={Stealth[length=1.5mm]},
  machine/.style={draw, rounded corners=0.8mm, fill=white},
  box/.style={draw, fill=white, minimum height=6mm, inner sep=1pt, align=center},
  kbox/.style={draw, fill=ln-box, minimum height=6mm, inner sep=1pt, align=center},
  nfs/.style={draw, fill=ln-accent!18, minimum height=6mm, inner sep=1pt, align=center},
  dev/.style={draw, fill=ln-accent!35, minimum height=5mm, inner sep=1pt, align=center},
  lab/.style={font=\scriptsize\sffamily, align=center},
  who/.style={font=\footnotesize\sffamily\bfseries, anchor=north west}]
  % ---------------- client machine
  \draw[machine] (0,0) rectangle (56,68);
  \node[who] at (1,67.5) {\iflangja{クライアントマシン}{client machine}};
  \node[box, minimum width=44mm] (app) at (28,56) {\iflangja{アプリケーション}{applications}};
  \draw[densely dashed, ln-muted] (1,50.5) -- (55,50.5);
  \node[kbox, minimum width=44mm] (sc) at (28,45) {\iflangja{システムコール層}{system call layer}};
  \node[kbox, minimum width=44mm] (vfs) at (28,35) {VFS};
  \node[kbox, minimum width=20mm] (lfs) at (16,24) {\iflangja{ローカル FS}{local FS}};
  \node[nfs, minimum width=20mm] (nc) at (40,24) {\iflangja{NFS クライアント}{NFS client}};
  \node[dev, minimum width=20mm] (d1) at (16,12) {\iflangja{ディスク}{disk}};
  \node[nfs, minimum width=20mm] (cs) at (40,12) {\iflangja{RPC スタブ}{RPC stub}};
  \draw[<->, thick] (app) -- (sc);
  \draw[<->, thick] (sc) -- (vfs);
  \draw[<->, thick] (lfs.north |- vfs.south) -- (lfs.north);
  \draw[<->, thick] (nc.north |- vfs.south) -- (nc.north);
  \draw[<->, thick] (lfs) -- (d1);
  \draw[<->, thick] (nc) -- (cs);
  % ---------------- server machine
  \draw[machine] (62,0) rectangle (118,68);
  \node[who] at (63,67.5) {\iflangja{サーバマシン}{server machine}};
  \node[nfs, minimum width=44mm] (ss) at (90,56) {\iflangja{RPC スタブ}{RPC stub}};
  \node[nfs, minimum width=44mm] (ns) at (90,45) {\iflangja{NFS サーバ}{NFS server}};
  \node[kbox, minimum width=44mm] (svfs) at (90,35) {VFS};
  \node[kbox, minimum width=44mm] (slfs) at (90,24) {\iflangja{ローカル FS}{local FS}};
  \node[dev, minimum width=44mm] (d2) at (90,12) {\iflangja{ディスク}{disk}};
  \draw[<->, thick] (ss) -- (ns);
  \draw[<->, thick] (ns) -- (svfs);
  \draw[<->, thick] (svfs) -- (slfs);
  \draw[<->, thick] (slfs) -- (d2);
  % ---------------- network
  \draw[<->, line width=1.4pt, ln-accent] (cs.east) -- (59,12) |- (ss.west);
  \node[lab, rotate=90, fill=white, inner sep=0.8pt, text=ln-accent] at (59,34) {\iflangja{ネットワーク（RPC）}{network (RPC)}};
  \node[lab, text=ln-muted, align=left, anchor=west] at (66,4.2)
    {\iflangja{要求はファイルハンドルで対象ファイルを指定}{requests name files by their file handles}};
\end{tikzpicture}

  \caption{NFS.  The client's VFS passes operations on NFS files to the NFS
    client, which sends them as RPC requests to the NFS server; the server
    performs them through its VFS and local file system.}
  \label{fig:l14-nfs}
\end{figure}

Which remote files a client can reach is determined when an NFS file
system is attached to the client's directory tree.  A server
\emph{exports} some of its directories, and a client makes an exported
directory available with the \term{mount}[マウント] operation, which
attaches it to a local directory, the \emph{mount point}.  Paths below the
mount point then name files on the server.  When a process opens such a
file, the NFS client looks up its path at the server, and the server
returns a file handle.

\begin{definition}[File handle]
  A \term{file handle}[ファイルハンドル] is the byte sequence by which the
  NFS client and the NFS server identify a file in their requests.  The
  server creates it, encoding the server-local identity of the file---the
  file system and the file within it---and the client treats it as opaque.
\end{definition}

The NFS client keeps the handle with its representation of the open file
and passes it in every subsequent request, such as a read, a write or a
request for the attributes of the file.  A server typically encodes in the
handle an identifier of the file system, the inode number of the file, and
a \emph{generation number} that changes whenever the inode is reused.  If
the file is removed on the server, by a local process or by another
client, requests with the handle fail with the error \emph{stale file
handle}: the handle refers to a file that no longer exists, and the
generation number prevents it from silently referring to a new file that
has reused the inode.

\begin{example}
  The directory \texttt{/home} of a client is mounted from a server; the
  mount operation has given the client the file handle of the exported
  directory.  The program below does not know this; its comments show the
  principal NFSv3 requests that the calls cause when nothing is cached.
  \begin{lstlisting}[language=C, numbers=none]
char buf[4096];
int fd = open("/home/ana/log.txt", O_RDONLY);
        /* LOOKUP "ana" in the handle of /home;
           LOOKUP "log.txt" in the handle of ana:
           the server returns the file handle   */
ssize_t n = read(fd, buf, sizeof buf);
        /* READ(file handle, offset 0, 4096)    */
close(fd);  /* read-only: no request is needed  */
\end{lstlisting}
\end{example}

\section{NFS versions}
\label{sec:l14-nfs-versions}

NFS has been in use since the 1980s and has gone through several
versions, of which NFSv3 \cite{rfc1813} and NFSv4 \cite{rfc7530} are in
current use.\source{P4L2, NFS versions; RFC~1813 \cite{rfc1813}; RFC~7530
\cite{rfc7530}; OSTEP ch.~49.}  They differ fundamentally in how they
treat state.  NFSv3 is a stateless protocol, although in practice it is
complemented by auxiliary protocols that keep state, such as a lock
manager.  NFSv4 is stateful: opening files, locking and the delegations
described below are part of the protocol.

\subsection{Caching}

For files that are not accessed concurrently, NFS provides session-based
semantics: a client writes its changes back to the server when it closes a
file, and checks for a newer version when it opens one.  In addition,
while a file is in use, the client periodically asks the server whether the
file has changed, by default every \SI{3}{\second} for files and every
\SI{30}{\second} for directories.\margin{On Linux these periods are the
mount options \texttt{acregmin} and \texttt{acdirmin}.}

\begin{example}
  Client~B last checked file $F$ with the server at
  $t=\SI{20.0}{\second}$, so its information about $F$ is considered
  current until \SI{23.0}{\second}.  Client~A closes $F$ at
  \SI{21.0}{\second}, writing its changes back.  A read by B at
  \SI{22.5}{\second} is served from B's cache and returns the old data; the
  first read by B after \SI{23.0}{\second} makes B check with the server,
  discover the change and fetch the new data.  A file that A creates at
  \SI{21.0}{\second} in a directory that B checked at \SI{20.0}{\second}
  may remain invisible to B until \SI{50.0}{\second}.
\end{example}

NFSv4 can avoid these checks.  The server may \emph{delegate} a file to a
client for a period of time.  While the delegation lasts, the client
opens the file and uses its cached data without checking with the server
for updates, and the server recalls the delegation when another client
needs access to the file that conflicts with it.

\begin{note}
  The resulting behaviour is neither UNIX semantics nor pure session
  semantics.  Between two checks a client may read data that another client
  has already changed at the server, which UNIX semantics excludes; and a
  periodic check can bring the changes of another client into a session
  that is still open, which session semantics excludes.
\end{note}

\subsection{Locking}

Locking in NFSv4 is based on leases.\margin{NFSv3 has no locking of its
own; clients lock files through the separate Network Lock Manager
protocol.}  The server grants a lock for a limited period, and before this
lease expires the client must either release the lock or renew the lease.
If a client fails or becomes unreachable, its lease runs out, and the
server can grant the lock to other clients without waiting indefinitely
for the failed one (\cref{fig:l14-lease}).  NFSv4 additionally supports
\emph{share reservations}: when a client opens a file, it declares the
access it requires---reading, writing or both---and the access it denies
to other clients.  Reservations that share reading but deny writing to
others act like a reader-writer lock\index{reader-writer lock}
(Lesson~10).

\begin{figure}[htbp]
  \centering
  % Lease-based locking: client 1 acquires a lock with a 30 s lease, renews it
% at 25 s and crashes at 40 s; the lease expires at 55 s and the server
% grants the lock to client 2, which has been waiting since 45 s.
% Usage: % Lease-based locking: client 1 acquires a lock with a 30 s lease, renews it
% at 25 s and crashes at 40 s; the lease expires at 55 s and the server
% grants the lock to client 2, which has been waiting since 45 s.
% Usage: % Lease-based locking: client 1 acquires a lock with a 30 s lease, renews it
% at 25 s and crashes at 40 s; the lease expires at 55 s and the server
% grants the lock to client 2, which has been waiting since 45 s.
% Usage: \input{figures/l14-lease}   (width ~116 mm)
\begin{tikzpicture}[x=1mm, y=1mm, font=\scriptsize\sffamily,
  >={Stealth[length=1.3mm]},
  row/.style={font=\scriptsize\sffamily, anchor=east},
  lab/.style={font=\scriptsize\sffamily, inner sep=0.8pt, align=center}]
  % time t (s) -> x = 18 + 1.4 t (mm)
  \node[row] at (16,16) {\iflangja{クライアント 1}{client 1}};
  \node[row] at (16,8) {\iflangja{サーバ}{server}};
  \node[row] at (16,0) {\iflangja{クライアント 2}{client 2}};
  \draw[ln-rule] (18,16) -- ({18+1.4*40},16);
  \draw[ln-rule, densely dotted] ({18+1.4*40},16) -- (116,16);
  \draw[ln-rule] (18,8) -- (116,8);
  \draw[ln-rule] (18,0) -- (116,0);
  % lock holder on the server
  \draw[fill=ln-box] ({18+1.4*0},6) rectangle ({18+1.4*55},10);
  \node[lab] at ({18+1.4*12},8) {\iflangja{クライアント 1 がロックを保持}{held by client 1}};
  \draw[fill=ln-accent!30] ({18+1.4*55},6) rectangle (116,10);
  \node[lab] at ({(18+1.4*55+116)/2},8) {\iflangja{クライアント 2}{client 2}};
  % lease expiry marks
  \foreach \t in {30,55} {
    \draw[densely dashed, ln-caution] ({18+1.4*\t},4.6) -- ({18+1.4*\t},11.4);
  }
  \node[lab, text=ln-caution, anchor=south] at ({18+1.4*55},11.4) {\iflangja{リース期限切れ}{lease expires}};
  \node[lab, text=ln-caution, anchor=north] at ({18+1.4*30},4.6) {\iflangja{当初の期限}{first expiry}};
  % client 1: acquire, renew, crash
  \draw[->, thick] ({18+1.4*0},14.8) -- ({18+1.4*0},10.3);
  \node[lab, anchor=south west] at ({18+1.4*0},17) {\iflangja{ロック獲得（リース 30 秒）}{lock (30 s lease)}};
  \draw[->, thick] ({18+1.4*25},14.8) -- ({18+1.4*25},10.3);
  \node[lab, anchor=south] at ({18+1.4*25},17) {\iflangja{更新}{renew}};
  \node[text=ln-caution, font=\normalsize\bfseries] at ({18+1.4*40},16) {$\times$};
  \node[lab, anchor=south, text=ln-caution] at ({18+1.4*40},17) {\iflangja{クラッシュ}{crash}};
  % client 2: request and grant
  \draw[->, thick] ({18+1.4*45},1.2) -- ({18+1.4*45},5.7);
  \node[lab, anchor=north] at ({18+1.4*45},-1.2) {\iflangja{ロック要求}{request}};
  \draw[densely dotted, thick] ({18+1.4*45},0) -- ({18+1.4*55},0);
  \draw[->, thick] ({18+1.4*55},5.7) -- ({18+1.4*55},1.2);
  \node[lab, anchor=north] at ({18+1.4*55},-1.2) {\iflangja{許可}{granted}};
  % time axis
  \draw[->] (18,-8) -- (117,-8);
  \foreach \t in {0,10,...,70} {
    \draw ({18+1.4*\t},-7.2) -- ({18+1.4*\t},-8.8);
    \node[lab, anchor=north] at ({18+1.4*\t},-9) {\t};
  }
  \node[lab, anchor=south east] at (117,-7.6) {\iflangja{時間（秒）}{time (s)}};
\end{tikzpicture}
   (width ~116 mm)
\begin{tikzpicture}[x=1mm, y=1mm, font=\scriptsize\sffamily,
  >={Stealth[length=1.3mm]},
  row/.style={font=\scriptsize\sffamily, anchor=east},
  lab/.style={font=\scriptsize\sffamily, inner sep=0.8pt, align=center}]
  % time t (s) -> x = 18 + 1.4 t (mm)
  \node[row] at (16,16) {\iflangja{クライアント 1}{client 1}};
  \node[row] at (16,8) {\iflangja{サーバ}{server}};
  \node[row] at (16,0) {\iflangja{クライアント 2}{client 2}};
  \draw[ln-rule] (18,16) -- ({18+1.4*40},16);
  \draw[ln-rule, densely dotted] ({18+1.4*40},16) -- (116,16);
  \draw[ln-rule] (18,8) -- (116,8);
  \draw[ln-rule] (18,0) -- (116,0);
  % lock holder on the server
  \draw[fill=ln-box] ({18+1.4*0},6) rectangle ({18+1.4*55},10);
  \node[lab] at ({18+1.4*12},8) {\iflangja{クライアント 1 がロックを保持}{held by client 1}};
  \draw[fill=ln-accent!30] ({18+1.4*55},6) rectangle (116,10);
  \node[lab] at ({(18+1.4*55+116)/2},8) {\iflangja{クライアント 2}{client 2}};
  % lease expiry marks
  \foreach \t in {30,55} {
    \draw[densely dashed, ln-caution] ({18+1.4*\t},4.6) -- ({18+1.4*\t},11.4);
  }
  \node[lab, text=ln-caution, anchor=south] at ({18+1.4*55},11.4) {\iflangja{リース期限切れ}{lease expires}};
  \node[lab, text=ln-caution, anchor=north] at ({18+1.4*30},4.6) {\iflangja{当初の期限}{first expiry}};
  % client 1: acquire, renew, crash
  \draw[->, thick] ({18+1.4*0},14.8) -- ({18+1.4*0},10.3);
  \node[lab, anchor=south west] at ({18+1.4*0},17) {\iflangja{ロック獲得（リース 30 秒）}{lock (30 s lease)}};
  \draw[->, thick] ({18+1.4*25},14.8) -- ({18+1.4*25},10.3);
  \node[lab, anchor=south] at ({18+1.4*25},17) {\iflangja{更新}{renew}};
  \node[text=ln-caution, font=\normalsize\bfseries] at ({18+1.4*40},16) {$\times$};
  \node[lab, anchor=south, text=ln-caution] at ({18+1.4*40},17) {\iflangja{クラッシュ}{crash}};
  % client 2: request and grant
  \draw[->, thick] ({18+1.4*45},1.2) -- ({18+1.4*45},5.7);
  \node[lab, anchor=north] at ({18+1.4*45},-1.2) {\iflangja{ロック要求}{request}};
  \draw[densely dotted, thick] ({18+1.4*45},0) -- ({18+1.4*55},0);
  \draw[->, thick] ({18+1.4*55},5.7) -- ({18+1.4*55},1.2);
  \node[lab, anchor=north] at ({18+1.4*55},-1.2) {\iflangja{許可}{granted}};
  % time axis
  \draw[->] (18,-8) -- (117,-8);
  \foreach \t in {0,10,...,70} {
    \draw ({18+1.4*\t},-7.2) -- ({18+1.4*\t},-8.8);
    \node[lab, anchor=north] at ({18+1.4*\t},-9) {\t};
  }
  \node[lab, anchor=south east] at (117,-7.6) {\iflangja{時間（秒）}{time (s)}};
\end{tikzpicture}
   (width ~116 mm)
\begin{tikzpicture}[x=1mm, y=1mm, font=\scriptsize\sffamily,
  >={Stealth[length=1.3mm]},
  row/.style={font=\scriptsize\sffamily, anchor=east},
  lab/.style={font=\scriptsize\sffamily, inner sep=0.8pt, align=center}]
  % time t (s) -> x = 18 + 1.4 t (mm)
  \node[row] at (16,16) {\iflangja{クライアント 1}{client 1}};
  \node[row] at (16,8) {\iflangja{サーバ}{server}};
  \node[row] at (16,0) {\iflangja{クライアント 2}{client 2}};
  \draw[ln-rule] (18,16) -- ({18+1.4*40},16);
  \draw[ln-rule, densely dotted] ({18+1.4*40},16) -- (116,16);
  \draw[ln-rule] (18,8) -- (116,8);
  \draw[ln-rule] (18,0) -- (116,0);
  % lock holder on the server
  \draw[fill=ln-box] ({18+1.4*0},6) rectangle ({18+1.4*55},10);
  \node[lab] at ({18+1.4*12},8) {\iflangja{クライアント 1 がロックを保持}{held by client 1}};
  \draw[fill=ln-accent!30] ({18+1.4*55},6) rectangle (116,10);
  \node[lab] at ({(18+1.4*55+116)/2},8) {\iflangja{クライアント 2}{client 2}};
  % lease expiry marks
  \foreach \t in {30,55} {
    \draw[densely dashed, ln-caution] ({18+1.4*\t},4.6) -- ({18+1.4*\t},11.4);
  }
  \node[lab, text=ln-caution, anchor=south] at ({18+1.4*55},11.4) {\iflangja{リース期限切れ}{lease expires}};
  \node[lab, text=ln-caution, anchor=north] at ({18+1.4*30},4.6) {\iflangja{当初の期限}{first expiry}};
  % client 1: acquire, renew, crash
  \draw[->, thick] ({18+1.4*0},14.8) -- ({18+1.4*0},10.3);
  \node[lab, anchor=south west] at ({18+1.4*0},17) {\iflangja{ロック獲得（リース 30 秒）}{lock (30 s lease)}};
  \draw[->, thick] ({18+1.4*25},14.8) -- ({18+1.4*25},10.3);
  \node[lab, anchor=south] at ({18+1.4*25},17) {\iflangja{更新}{renew}};
  \node[text=ln-caution, font=\normalsize\bfseries] at ({18+1.4*40},16) {$\times$};
  \node[lab, anchor=south, text=ln-caution] at ({18+1.4*40},17) {\iflangja{クラッシュ}{crash}};
  % client 2: request and grant
  \draw[->, thick] ({18+1.4*45},1.2) -- ({18+1.4*45},5.7);
  \node[lab, anchor=north] at ({18+1.4*45},-1.2) {\iflangja{ロック要求}{request}};
  \draw[densely dotted, thick] ({18+1.4*45},0) -- ({18+1.4*55},0);
  \draw[->, thick] ({18+1.4*55},5.7) -- ({18+1.4*55},1.2);
  \node[lab, anchor=north] at ({18+1.4*55},-1.2) {\iflangja{許可}{granted}};
  % time axis
  \draw[->] (18,-8) -- (117,-8);
  \foreach \t in {0,10,...,70} {
    \draw ({18+1.4*\t},-7.2) -- ({18+1.4*\t},-8.8);
    \node[lab, anchor=north] at ({18+1.4*\t},-9) {\t};
  }
  \node[lab, anchor=south east] at (117,-7.6) {\iflangja{時間（秒）}{time (s)}};
\end{tikzpicture}

  \caption{Lease-based locking.  Client 1 renews its lease at
    \SI{25}{\second} and crashes at \SI{40}{\second}; when the lease
    expires at \SI{55}{\second}, the server grants the lock to client 2.}
  \label{fig:l14-lease}
\end{figure}

\begin{example}
  A server grants leases of \SI{30}{\second}.  Client~1 acquires a lock at
  $t=0$ and renews the lease at \SI{25}{\second}, extending it to
  \SI{55}{\second}; it crashes at \SI{40}{\second}.  Client~2 requests the
  lock at \SI{45}{\second} and waits.  At \SI{55}{\second} the lease of
  client~1 expires without renewal, and the server grants the lock to
  client~2 (\cref{fig:l14-lease}).  Had client~1 not crashed but merely lost
  contact with the server, it would have to stop using the lock at
  \SI{55}{\second}, because it could not know whether the lock had been
  granted to another client.
\end{example}

\begin{keypoint}
  NFS clients reach remote files through the VFS, the NFS client and RPC,
  and identify files by server-generated file handles.  NFSv3 is
  stateless and NFSv4 stateful.  NFS combines session-based caching with
  periodic checks, by default every \SI{3}{\second} for files and every
  \SI{30}{\second} for directories; NFSv4 adds delegations, lease-based
  locking and share reservations.
\end{keypoint}

\section{The Sprite distributed file system}
\label{sec:l14-sprite}

\term{Sprite}[Sprite] was a distributed operating system developed at the
University of California, Berkeley, in the 1980s, whose file system let
the workstations of a network share files stored on servers.\source{P4L2,
Sprite distributed file system, Sprite DFS access pattern analysis, Sprite
DFS from analysis to design; Nelson et al.\ \cite{nelson1988}.}  Nelson,
Welch and Ousterhout \cite{nelson1988} describe how its file system caches
data.  The paper is instructive above all for its design process: each
design decision is derived from, and justified by, an analysis of how
files are actually accessed.  \Cref{tab:l14-sprite-analysis} lists the
observations and the conclusions drawn from them.

\begin{table}[htbp]
  \centering
  \caption{Access patterns and their consequences for the design of
    Sprite.}
  \label{tab:l14-sprite-analysis}
  \small
  \begin{tabular}{@{}>{\raggedright\arraybackslash}p{0.4\linewidth}>{\raggedright\arraybackslash}p{0.54\linewidth}@{}}
    \toprule
    Observation & Consequence \\
    \midrule
    \SI{33}{\percent} of all file accesses are writes & Caching pays off
      for the majority of accesses, which are reads; but propagating every
      write at once, as UNIX semantics requires, is too expensive. \\
    \SI{75}{\percent} of files are open for less than \SI{0.5}{\second},
      \SI{90}{\percent} for less than \SI{10}{\second} & Session semantics
      still has too much overhead: with such short sessions, updates at
      every open and close would reach the server very frequently. \\
    \SIrange{20}{30}{\percent} of new data are deleted within
      \SI{30}{\second}, \SI{50}{\percent} within \SI{5}{\minute} & Writing
      data back at close is not necessary; much of the data are deleted
      before any other client needs them. \\
    Concurrent access to a file is rare & Concurrent access need not be
      optimized, but it must be supported. \\
    \bottomrule
  \end{tabular}
\end{table}

From these observations the design of Sprite follows:
\begin{itemize}
  \item Clients cache file blocks and write modified blocks back to the
    server with a delay of 30 seconds, as described below.
  \item When a client opens a file that another client has been writing,
    the server first collects the dirty blocks from that client.
  \item Every open goes to the server, and directories are not cached at
    the clients.
  \item When a file is written concurrently, caching of that file is
    disabled and all its operations go to the server.
\end{itemize}

The first point defines the write policy.  Keeping modified blocks in the
client cache and writing them back only after a delay, rather than on
every write or close, is called \term{delayed write}[遅延書き込み] or
delayed write-back.  In Sprite, every 30 seconds a client writes back the
dirty blocks that have not been modified during the last 30 seconds.
Blocks that are still being modified are therefore not transferred
repeatedly, and data deleted within 30 seconds of being written are never
transferred at all.  A block reaches the server between 30 and 60 seconds
after its last modification.

\begin{example}
  A client writes back at $t = 0, 30, 60, 90, \ldots$~seconds.  Block
  $b_1$ is modified at \SI{12}{\second}; at \SI{30}{\second} it has been
  unmodified for only \SI{18}{\second}, so it is written back at
  \SI{60}{\second}.  Block $b_2$ is modified at \SI{12}{\second} and again
  at \SI{50}{\second}; at \SI{60}{\second} it has been unmodified for
  \SI{10}{\second}, so it is written back at \SI{90}{\second}.  Block
  $b_3$ belongs to a temporary file that is written at \SI{20}{\second}
  and deleted at \SI{45}{\second}; it is never sent to the server.
\end{example}

The semantics Sprite provides depend on how a file is shared.

\begin{definition}[Sequential and concurrent write sharing]
  A file is subject to
  \term{sequential write sharing}[逐次書き込み共有] when one client
  modifies it and closes it and another client opens it later: the clients
  use the file one after another.  It is subject to
  \term{concurrent write sharing}[並行書き込み共有] when it is open on
  several clients at the same time and at least one of them has it open for
  writing.
\end{definition}

Under sequential write sharing Sprite caches the file, and every client
that opens it sees the data written by the previous writer.  Under
concurrent write sharing Sprite does not cache the file; the server
executes all reads and writes in order, which yields UNIX semantics.

\section{File access operations in Sprite}
\label{sec:l14-sprite-ops}

Consider a file that several clients read and one or more clients
write.\source{P4L2, file access operations in Sprite; Nelson et al.\
\cite{nelson1988}.}  All opens go through the server, all clients cache
blocks while the file is cacheable, and a writer keeps a timestamp for
every block it modifies, as delayed write requires.  Clients and server
keep the following state for each file (\cref{fig:l14-sprite}):
\begin{description}
  \item[Client] whether caching is enabled for the file; the cached blocks;
    the time of the last modification of each dirty block; and the version
    number of the cached data.
  \item[Server] the clients that have the file open for reading and for
    writing; the client that wrote the file last; the current version
    number; and whether the file is cacheable.
\end{description}

\begin{figure}[htbp]
  \centering
  % Sprite: per-file state on clients and server, and the messages when W3
% opens a file that W2 has open for writing (concurrent write sharing).
% Usage: % Sprite: per-file state on clients and server, and the messages when W3
% opens a file that W2 has open for writing (concurrent write sharing).
% Usage: % Sprite: per-file state on clients and server, and the messages when W3
% opens a file that W2 has open for writing (concurrent write sharing).
% Usage: \input{figures/l14-sprite}   (width ~118 mm)
\begin{tikzpicture}[x=1mm, y=1mm, font=\scriptsize\sffamily,
  >={Stealth[length=1.5mm]},
  machine/.style={draw, rounded corners=0.8mm},
  st/.style={font=\scriptsize\sffamily, anchor=north west, inner sep=0pt},
  lab/.style={font=\scriptsize\sffamily, align=center, inner sep=0.8pt},
  who/.style={font=\scriptsize\sffamily\bfseries, anchor=north west}]
  % client W2 (writer)
  \draw[machine, fill=white] (0,27) rectangle (46,47);
  \node[who] at (1,46.5) {\iflangja{クライアント $W_2$（書き込み中）}{client $W_2$ (writing)}};
  \node[st] at (2,41.5) {\begin{tabular}{@{}l@{\quad}l@{}}
      \iflangja{キャッシュ}{caching} & \iflangja{有効 → 無効}{on $\to$ off}\\
      \iflangja{ブロック}{blocks} & 0, 1, 5\\
      \iflangja{ダーティ}{dirty} & 1 (41\,s), 5 (44\,s)\\
      \iflangja{版}{version} & 6
    \end{tabular}};
  % client W3 (opening)
  \draw[machine, fill=white] (72,27) rectangle (118,47);
  \node[who] at (73,46.5) {\iflangja{クライアント $W_3$（open 中）}{client $W_3$ (opening)}};
  \node[st] at (74,41.5) {\begin{tabular}{@{}l@{\quad}l@{}}
      \iflangja{キャッシュ}{caching} & \iflangja{無効}{off}\\
      \iflangja{ブロック}{blocks} & ---\\
      \iflangja{版}{version} & 7
    \end{tabular}};
  % server
  \draw[machine, fill=ln-box] (33,-8) rectangle (85,13);
  \node[who] at (34,12.5) {\iflangja{サーバ}{server}};
  \node[st] at (35,7.5) {\begin{tabular}{@{}l@{\quad}l@{}}
      \iflangja{読み手}{readers} & ---\\
      \iflangja{書き手}{writers} & $W_2$, $W_3$\\
      \iflangja{版}{version} & $6 \to 7$\\
      \iflangja{キャッシュ可能}{cacheable} & \iflangja{可 → 不可}{yes $\to$ no}
    \end{tabular}};
  % (1) open by W3
  \draw[->, thick] (108,27) |- (85,2);
  \node[lab, anchor=south] at (96,2.6) {\iflangja{(1) open\\（書き込み）}{(1) open\\for writing}};
  % (2) dirty blocks from W2
  \draw[<->, thick, ln-accent] (33,2) -| (12,27);
  \node[lab, anchor=west, text=ln-accent] at (13,16.5) {\iflangja{(2) ダーティ\\ブロック}{(2) dirty\\blocks}};
  % (3) disable caching
  \draw[->, densely dashed, ln-caution, thick] (59,13) -- (59,20) -| (36,27);
  \draw[->, densely dashed, ln-caution, thick] (59,20) -| (82,27);
  \node[lab, anchor=south, text=ln-caution, fill=white] at (59,20.6) {\iflangja{(3) キャッシュ無効化}{(3) disable caching}};
\end{tikzpicture}
   (width ~118 mm)
\begin{tikzpicture}[x=1mm, y=1mm, font=\scriptsize\sffamily,
  >={Stealth[length=1.5mm]},
  machine/.style={draw, rounded corners=0.8mm},
  st/.style={font=\scriptsize\sffamily, anchor=north west, inner sep=0pt},
  lab/.style={font=\scriptsize\sffamily, align=center, inner sep=0.8pt},
  who/.style={font=\scriptsize\sffamily\bfseries, anchor=north west}]
  % client W2 (writer)
  \draw[machine, fill=white] (0,27) rectangle (46,47);
  \node[who] at (1,46.5) {\iflangja{クライアント $W_2$（書き込み中）}{client $W_2$ (writing)}};
  \node[st] at (2,41.5) {\begin{tabular}{@{}l@{\quad}l@{}}
      \iflangja{キャッシュ}{caching} & \iflangja{有効 → 無効}{on $\to$ off}\\
      \iflangja{ブロック}{blocks} & 0, 1, 5\\
      \iflangja{ダーティ}{dirty} & 1 (41\,s), 5 (44\,s)\\
      \iflangja{版}{version} & 6
    \end{tabular}};
  % client W3 (opening)
  \draw[machine, fill=white] (72,27) rectangle (118,47);
  \node[who] at (73,46.5) {\iflangja{クライアント $W_3$（open 中）}{client $W_3$ (opening)}};
  \node[st] at (74,41.5) {\begin{tabular}{@{}l@{\quad}l@{}}
      \iflangja{キャッシュ}{caching} & \iflangja{無効}{off}\\
      \iflangja{ブロック}{blocks} & ---\\
      \iflangja{版}{version} & 7
    \end{tabular}};
  % server
  \draw[machine, fill=ln-box] (33,-8) rectangle (85,13);
  \node[who] at (34,12.5) {\iflangja{サーバ}{server}};
  \node[st] at (35,7.5) {\begin{tabular}{@{}l@{\quad}l@{}}
      \iflangja{読み手}{readers} & ---\\
      \iflangja{書き手}{writers} & $W_2$, $W_3$\\
      \iflangja{版}{version} & $6 \to 7$\\
      \iflangja{キャッシュ可能}{cacheable} & \iflangja{可 → 不可}{yes $\to$ no}
    \end{tabular}};
  % (1) open by W3
  \draw[->, thick] (108,27) |- (85,2);
  \node[lab, anchor=south] at (96,2.6) {\iflangja{(1) open\\（書き込み）}{(1) open\\for writing}};
  % (2) dirty blocks from W2
  \draw[<->, thick, ln-accent] (33,2) -| (12,27);
  \node[lab, anchor=west, text=ln-accent] at (13,16.5) {\iflangja{(2) ダーティ\\ブロック}{(2) dirty\\blocks}};
  % (3) disable caching
  \draw[->, densely dashed, ln-caution, thick] (59,13) -- (59,20) -| (36,27);
  \draw[->, densely dashed, ln-caution, thick] (59,20) -| (82,27);
  \node[lab, anchor=south, text=ln-caution, fill=white] at (59,20.6) {\iflangja{(3) キャッシュ無効化}{(3) disable caching}};
\end{tikzpicture}
   (width ~118 mm)
\begin{tikzpicture}[x=1mm, y=1mm, font=\scriptsize\sffamily,
  >={Stealth[length=1.5mm]},
  machine/.style={draw, rounded corners=0.8mm},
  st/.style={font=\scriptsize\sffamily, anchor=north west, inner sep=0pt},
  lab/.style={font=\scriptsize\sffamily, align=center, inner sep=0.8pt},
  who/.style={font=\scriptsize\sffamily\bfseries, anchor=north west}]
  % client W2 (writer)
  \draw[machine, fill=white] (0,27) rectangle (46,47);
  \node[who] at (1,46.5) {\iflangja{クライアント $W_2$（書き込み中）}{client $W_2$ (writing)}};
  \node[st] at (2,41.5) {\begin{tabular}{@{}l@{\quad}l@{}}
      \iflangja{キャッシュ}{caching} & \iflangja{有効 → 無効}{on $\to$ off}\\
      \iflangja{ブロック}{blocks} & 0, 1, 5\\
      \iflangja{ダーティ}{dirty} & 1 (41\,s), 5 (44\,s)\\
      \iflangja{版}{version} & 6
    \end{tabular}};
  % client W3 (opening)
  \draw[machine, fill=white] (72,27) rectangle (118,47);
  \node[who] at (73,46.5) {\iflangja{クライアント $W_3$（open 中）}{client $W_3$ (opening)}};
  \node[st] at (74,41.5) {\begin{tabular}{@{}l@{\quad}l@{}}
      \iflangja{キャッシュ}{caching} & \iflangja{無効}{off}\\
      \iflangja{ブロック}{blocks} & ---\\
      \iflangja{版}{version} & 7
    \end{tabular}};
  % server
  \draw[machine, fill=ln-box] (33,-8) rectangle (85,13);
  \node[who] at (34,12.5) {\iflangja{サーバ}{server}};
  \node[st] at (35,7.5) {\begin{tabular}{@{}l@{\quad}l@{}}
      \iflangja{読み手}{readers} & ---\\
      \iflangja{書き手}{writers} & $W_2$, $W_3$\\
      \iflangja{版}{version} & $6 \to 7$\\
      \iflangja{キャッシュ可能}{cacheable} & \iflangja{可 → 不可}{yes $\to$ no}
    \end{tabular}};
  % (1) open by W3
  \draw[->, thick] (108,27) |- (85,2);
  \node[lab, anchor=south] at (96,2.6) {\iflangja{(1) open\\（書き込み）}{(1) open\\for writing}};
  % (2) dirty blocks from W2
  \draw[<->, thick, ln-accent] (33,2) -| (12,27);
  \node[lab, anchor=west, text=ln-accent] at (13,16.5) {\iflangja{(2) ダーティ\\ブロック}{(2) dirty\\blocks}};
  % (3) disable caching
  \draw[->, densely dashed, ln-caution, thick] (59,13) -- (59,20) -| (36,27);
  \draw[->, densely dashed, ln-caution, thick] (59,20) -| (82,27);
  \node[lab, anchor=south, text=ln-caution, fill=white] at (59,20.6) {\iflangja{(3) キャッシュ無効化}{(3) disable caching}};
\end{tikzpicture}

  \caption{Per-file state in Sprite.  When $W_3$ opens the file that $W_2$
    has open for writing (1), the server collects $W_2$'s dirty blocks (2)
    and disables caching on both clients (3).}
  \label{fig:l14-sprite}
\end{figure}

\paragraph{Sequential write sharing.}  Each time a client opens the file
for writing, the server increments the version number.  A client that
opens the file compares the version number returned by the server with
that of its cached blocks, and discards the blocks if they are older.
Because of delayed write, the client that wrote the file last may still
hold dirty blocks of it even after closing it.  The server therefore asks
that client to write them back before another client opens the file, so
that the new client sees the latest data.  Caching remains enabled.

\paragraph{Concurrent write sharing.}  When a client opens the file for
writing while another client has it open, or for reading while another
client has it open for writing, the server again collects the dirty blocks
of the writer, marks the file as not cacheable, and notifies all clients
that have it open.  From then on no client caches the file, and every read
and write goes to the server.  The file becomes cacheable again once all
clients have closed it.

\begin{example}
  \Cref{tab:l14-sprite-trace} traces the server state for file $F$,
  initially at version~4, as clients $R_1$, $W_1$, $W_2$ and $W_3$ open and
  close it.  In step~5 $W_1$ has closed $F$ but still holds dirty blocks,
  which the server collects before $W_2$ proceeds; the sharing is
  sequential, and $W_2$ may cache $F$.  In step~6 $W_3$ opens $F$ while
  $W_2$ still has it open for writing (\cref{fig:l14-sprite}), and caching
  is disabled.  In step~7 $R_1$ finds that its cached blocks of version~4
  are outdated; it discards them and reads from the server.
\end{example}

\begin{table}[htbp]
  \centering
  \caption{Server state for file $F$ in Sprite.  Version and
    cacheability are shown after each step.}
  \label{tab:l14-sprite-trace}
  \small
  \begin{tabular}{@{}r>{\raggedright\arraybackslash}p{0.26\linewidth}>{\raggedright\arraybackslash}p{0.3\linewidth}llc@{}}
    \toprule
    & Event & Server action & Open by & Ver. & Cache \\
    \midrule
    1 & $R_1$ opens for reading & --- & $R_1$ & 4 & yes \\
    2 & $R_1$ closes & --- & --- & 4 & yes \\
    3 & $W_1$ opens for writing & increments version & $W_1$ & 5 & yes \\
    4 & $W_1$ writes, closes & --- (dirty blocks stay at $W_1$) & --- & 5
      & yes \\
    5 & $W_2$ opens for writing & collects $W_1$'s dirty blocks; increments
      version & $W_2$ & 6 & yes \\
    6 & $W_3$ opens for writing & collects $W_2$'s dirty blocks; increments
      version; disables caching & $W_2$, $W_3$ & 7 & no \\
    7 & $R_1$ opens for reading & tells $R_1$ not to cache & $R_1$, $W_2$,
      $W_3$ & 7 & no \\
    8 & $W_2$, $W_3$, $R_1$ close & enables caching & --- & 7 & yes \\
    \bottomrule
  \end{tabular}
\end{table}

\begin{keypoint}[Lesson summary]
  A distributed file system serves files from networked machines through
  the ordinary file system interface, with one server, replicated or
  partitioned servers, or peers.  Practical designs lie between the
  upload/download and remote access models: clients cache blocks and
  interact frequently with a stateful server, which must keep the cached
  copies coherent.  File sharing semantics---UNIX, session, periodic
  updates with leases, immutable files, transactions---define what clients
  observe, and may differ for files and directories.  Replication improves
  availability and load balance at the cost of complex writes;
  partitioning scales with the size of the file system.  NFS uses file
  handles over RPC, stateless NFSv3 or stateful NFSv4, session-based
  caching with periodic checks, and lease-based locking.  Sprite derived
  its design from measured access patterns: delayed write with a
  30-second delay, caching under sequential write sharing, and no caching
  under concurrent write sharing.
\end{keypoint}

\end{document}
